% detect interpreter: pdflatex or latex
\newif\ifpdf
\ifx\pdfoutput\undefined
\pdffalse
\else
\pdftrue
\fi

\documentclass[10pt]{book}

\usepackage{graphicx} % enhanced graphics usage
\usepackage{a4} % A4 paperformat
\usepackage{isolatin1} % for input of 8 bits character
\usepackage{makeidx} % enable indexing
\usepackage{verbatim} % better verbatim environment
\usepackage{epsfig}
\usepackage{xspace} % add extra space at the end of the word if necessary
\usepackage{color} % provides standard LaTeX colors
\usepackage{fancyhdr}
\usepackage{float} % float environment enhancements
\usepackage{alltt} % defines alltt environment which is like verbatim, but allows some additional formating
\usepackage{listings} % pritty print of code-listings
\ifpdf
% this package should be loaded last
\usepackage[colorlinks=true,
            pdfstartview=FitV,
            linkcolor=blue,
            citecolor=blue,
            urlcolor=blue]{hyperref}
\fi

\graphicspath{{../fig/}}

%% insert the common definitions
\usepackage{a4} % A4 paperformat
\usepackage{isolatin1} % for input of 8 bits character
\usepackage{makeidx} % enable indexing
\usepackage{verbatim} % better verbatim environment
\usepackage{epsfig}
\usepackage{xspace} % add extra space at the end of the word if necessary
\usepackage{color} % provides standard LaTeX colors
\usepackage{graphicx} % enhanced graphics usage
\usepackage{fancyhdr}
\usepackage{float} % float environment enhancements
\usepackage{alltt} % defines alltt environment which is like verbatim, but allows some additional formating
\usepackage{listings} % pritty print of code-listings

%\usepackage[german]{babel} % multilingual package
%\usepackage[normal]{subfigure} % support for the inclusion of small subfigures
%\usepackage{latexsym} % some extra mathematical symbols
%\usepackage{exscale} % implements scaling of the `cmex' fonts


% configure the listings-environment
\lstloadlanguages{C++}
\lstset{
  basicstyle=\sf,
  commentstyle=\hfill \sl,
  language=C++,
  defaultdialect=ANSI,
  flexiblecolumns=true,
  indent=\parindent,
  extendedchars
  %%  texcl
  %% formfeed={}
}


% command definitions

% logo of the miro-software writing
\newcommand{\miro}{\textit{Miro}\xspace}
% logo of the sparrow writing
\newcommand{\sparrow}{\textit{Sparrow-99}\xspace}
% logo of the b21 writing
\newcommand{\bXXI}{\textit{B21}\xspace}


% font definitions
\newcommand{\Xbombastic}{\fontfamily{ppl}\fontseries{b}\fontsize{1.5in}{1.8in}\selectfont}


% color definitions
\definecolor{LightGrey}{rgb}{0.9,0.9,0.9}


% page definitions
\setlength{\parindent}{0pt} % Keine Absatzeinr�ckung
\setlength{\parskip}{7pt plus 1pt minus 1pt} % Absatz Abstand 7pt


% misc
\makeindex
\bibliographystyle{plain}

%=============================== Examples ====================================

%% code-parts as an extra paragraph
%% \begin{lstlisting}[frame=tb]{}
%% for (unsigned int i = 0; i < sonarScan.length(); ++i)
%%   cout << sonarScan.[i] << " ";
%% \end{lstlisting}

%% Included code-parts from a file
%% \lstinputlisting[frame=tb]{simpleTestClient.cc}

%=============================== Examples (end) ====================================

%%% Local Variables:  
%%% mode: latex
%%% TeX-master: "miro_manual"
%%% End: 


%% special definitions
\setcounter{tocdepth}{3}

%% =============================================================================

\begin{document}
%\maketitle

\thispagestyle{empty}
\begin{center}
  \vfill
  \includegraphics[width=5in]{signature}\\
  \vspace{30 mm}
  {\Xbombastic \miro}\\
  \vspace{10 mm}
  {\LARGE \textbf{Manual}}\\
  \vspace{10 mm}
  \textbf{Version 0.9.4}\\
  \vspace{10 mm}
  \today\\
  \vfill
\end{center}

\newpage

\begin{center}
  \includegraphics[width=5in]{pincolla}
  \bigskip
  For more paintings, see {\tt http://www.bcn.fjmiro.es/}
\end{center}

\pagestyle{headings}

\newpage
\tableofcontents

%% =============================================================================

\chapter{Introduction}

\miro is a distributed object oriented framework for mobile robot
control, based on CORBA (Common Object Request Broker
Architecture) \cite{OMG:CORBA-standart} technology. The \miro core
components have been developed in C++ for Linux. But due to the
programming language independency of CORBA further components can be
written in any language and on any platform that provides CORBA
implementations.

The \miro core components have been developed under the aid of ACE
(Adaptive Communications Environment) \cite{ACE-Paper, ACE-Manual}, an
object oriented multi-platform framework for OS-independent
interprocess, network and real time communication. They use TAO (The
ACE ORB) \cite{TAO-paper} as their ORB (Object Request Broker), a
CORBA implementation designed for high performance and real time
applications. Therefore \miro should be easily portable to any other
OS, where ACE and TAO run on.  These are many Unix clones, Windows NT
and some real time operating systems.

\miro is currently available for the RWI B21 platform (pre rflex), the
Activmedia Pioneer family (namely the Performance PeopleBot), and the
vi\sparrow architecture developed at the University of Ulm.  We are
convinced, that other ports can be done straight forward.

\miro was built since the existing robot control architectures didn't
suffice our needs of usability, reliability, scalability and
portability.

We used C++ due to its advantages in big projects, since it was
especially designed for big projects. We have learned the hard way,
that this is a serious requirement, for projects like a mobile robot
control architecture.

We use multi processing, multithreading and the CORBA technology since
most robotics applications are inherently concurrent and distributed.
The hardware devices, like sensors and actuators run concurrently and
due to the constant lack of computing power especially in computer
vision, tend to reside on multiple computers. And as soon as
cooperative behavior of multiple autonomous robots is of interest, the
system as a whole is distributed anyways.

%We used an object oriented approach, since real live objects are
%concerned. And on the other hand, OO is just state of the art and you
%need good arguments to start a new project ignoring object
%orientation.

And last but not least we use ACE and TAO since these are multi
platform, high performance libraries which proved to be very
sophisticated in terms of usability, portability and scalability. ---
Additionally, they are open source libraries. They haven't
disappointed us yet, to the contrary.

\section{The \miro Group}

The \miro core developers are (in alphabetical order):
\begin{itemize}
  \item Stefan Enderle
  \item Gerhard Kraetzschmar
  \item Gerd Mayer
  \item Guillem Pages
  \item Stefan Sablatn\"og
  \item Steffen Simon
  \item Hans Utz
\end{itemize}

\section{\miro Directory Structure}

\miro comes with the full source code and documentation as well as a
set of test programs and examples that should facilitate your first
steps when writing your own programs. To help you navigate through the
directory tree of a \miro distribution, we give a brief overview of
the directories present in the \miro root directory:

\begin{description}
\item[bin:]
  This directory contains links to the binaries of \miro.
  For an explanation of the individual binaries, see Chapter
  \ref{sec:services} about the robot services.
\item[doc:]
  Here, the available documentation is gathered. In the
  subdirectory {\tt tex} you find the postscript version of this
  manual and in the directory {\tt html} resides the auto generated
  online documentation of all \miro classes and their methods
  (the starting page of this documentation resides at: \\
  {\tt \$MIRO\_ROOT/doc/html/idl/index.html} respectively \\
  {\tt \$MIRO\_ROOT/doc/html/cpp/index.html}).
\item[etc:]
  Config files for the individual robots.
\item[examples:]
  Examples on how to use individual interfaces of \miro. If you want
  to write your own programs utilizing the \miro framework, this is a
  good place to look for inspirations.
\item[idl:]
  The IDL sources of the CORBA interfaces and data types.
\item[performance-tests:]
  Some tests that measure the performance of \miro's
  services. --- Not too much there at the moment.
\item[scripts:]
  Utilities for source code formatting and handling sequences of
  datafiles.
\item[src:]
  Here, all sources of the \miro services reside. New services for
  further robot platforms should go in here.
\item[templates:]
  Templates for Makefiles and the headers for source files. Copy the
  corresponding template, if you want to start a new subproject within
  the \miro source hierarchy. If you start your own new project on top
  of \miro, the Makefile templates might still be useful for you. They
  handle all the stuff concerning ACE/TAO and the multi-platform build
  process. See Chapter \ref{sec:makefile} for details.
\item[tests:]
  Small test programs to monitor or test isolated interfaces of the
  \miro robot control architecture.
\item[utils:]
  Utilities made for convenience. No magic to be expected here.
\end{description}

%=============================== Tutorial ====================================

\chapter{Definitions}


\section{Units}

When it comes to parameter passing, distances are specified in mm and
angles in radiant. Velocities are specified in mm/s and rad/s
respectively. 

\section{Coordinates}

\miro uses the cartesian coordinate space. Angles are specified
mathematically, that is, counter-clock wise.

%mm, rad


\begin{figure}[!ht]
  \begin{center}
    \ifpdf
    \includegraphics[width=8cm]{../fig/grid_coords.png}
    \else
    \includegraphics[width=8cm]{../fig/grid_coords.eps}
    \fi
  \end{center}
\end{figure}


%%% Local Variables: 
%%% mode: latex
%%% TeX-master: "miro_manual"
%%% End: 

\chapter{Using \miro}
\label{sec:using}

In order to use \miro programs, the following steps have to be performed:
\begin{itemize}
  \item start the naming service
  \item start the services
  \item start the client program(s)
\end{itemize}

These steps are described in the following.


\section{Starting the CORBA Naming Service}

The CORBA naming service is some sort of telephone book service for
the lookup of remote objects. Its organized like a file system. All
entries in the naming service can be either object references, the
pointer to a remote object, or naming context references that point to
another subdirectory of the naming service.  All CORBA objects of
\miro are registered by {\em name} at a CORBA {\em naming service}. In
the \miro project, the TAO implementation of the CORBA naming service
is used and therefore must be started before any service is run.

The TAO naming service is started by the following command:

\begin{lstlisting}[frame=tb]{}
cd ${TAO_ROOT}/TAO/orbsvcs/Naming_Service
./Naming_Service
\end{lstlisting}

In order to get debug messages it can also be started like this:

\begin{lstlisting}[frame=tb]{}
cd ${TAO_ROOT}/TAO/orbsvcs/Naming_Service
./Naming_Service -ORBDebug -ORBDebugLevel <level>
\end{lstlisting}

where {\tt <level>} is a debug level of 1 to 10 (10 gives most debug information).


\subsection{Naming Service Lookup by IP-Multicast}

Programs that use the TAO ORB can locate the naming service by IP
multicast as long as it resided in the same subnet, or within the same
multicast group. This is a special (proprietary) feature of the TAO
naming service. You have to enable this
feature explicitly by the command line option {\tt -m 1}. To avoid
conflicts, there should only be one naming service in a subnet. To
make sure no other naming service is active or to check the entries in
an active naming service, TAO provides a utility called {\tt nslist}.
It can be started by calling:

\begin{lstlisting}[frame=tb]{}
${TAO_ROOT}/TAO/utils/nslist/nslist 
\end{lstlisting}
%%$

It lists the contents of the currently running naming service and
returns an error, if no naming service is found.

\subsection{Naming Service Lookup by Environment Variable}

Additionally it is possible to not use the multicast feature of the
TAO naming service and provide the IOR of the naming service
as command line arguments or via environment variables to any program
that uses TAO. This is explained in detail in the TAO documentation
under:

\begin{lstlisting}[frame=tb]{}
${TAO_ROOT}/TAO/docs/INS.html
\end{lstlisting}
%% $

Note, that for retrieving the IOR of an already running naming service the
little tool {\tt nsIOR} resides in the bin directory of \miro.

You can also tell the Naming Service to listen to a specific endpoint
instead of choosing its own by the use of the {\tt -ORBEndpoint}
command line option. Ensuring that the IOR of the naming service is
the same, after each restart. 

\begin{lstlisting}[frame=tb]{}
\${MIRO_ROOT}/scritps/naming_service
\end{lstlisting}

is a small init script for the SuSE Linux distribution, that can be
installed in {\tt /etc/rc.d} for starting the naming service
automatically at boot time.


\subsection{Naming Contexts}

To provide the possibility to operate multiple robots in parallel,
\miro uses the concept of naming contexts to distinguish
between services of individual robots. One can imagine {\em naming
  contexts} as different folders, entries in different folders do not
interact. The default naming context \miro services use is
{\tt Miro}. The canonical naming convention in a multi robot
scenario is, that every robot uses its own name. Therefore every
\miro service and client accepts an additional command-line argument
that sets the {\em naming context}: 
{\tt -MiroNamingContext <name>}. Or as short form: 
{\tt -MNC <name>}.
To access services of different robots simultaneously, the
{\em naming context} has to be specified for the resolution of
each individual service (see Chapter \ref{sec:clients}).


\section{Starting Services}

When the naming service is running and can be accessed from all
machines, the most difficult part of the setup of the distributed
system environment is mastered. From now on we are back to the
unbounded problem set of robotics.

Simply change to the directory {\tt \$MIRO\_ROOT/bin} and start the
service you want to run. For example {\tt b21Base},
{\tt laserServer}, {\tt videoService}, {\tt dpPanTilt}. For
an extensive list of all available services see Chapter
\ref{sec:services}. Since the services depend on the {\tt *.conf}
files provided in this directory, strange effects are bound to occur
if the services are started from a different directory. Keep in mind,
that the hardware dependent services must run on the computers the
corresponding hardware is attached to.

%Services register their interfaces at the naming service. To allow
%multiple robots to register the same interfaces, each robot chooses a
%naming context in the naming service. This can be seen as
%subdirectories within a file system. Files with the same name can live
%peacefully side by side within different directories. To chose a naming
%context different than the default (which is {\tt Miro}) the command
%line option {\tt -MiroNamingContext context} can be used.

When a service is ended disgracefully, it probably doesn't deregister
its interfaces at the naming service. This will prevent the restated
service from registering the new interface references. TAO provides a
utility to manually remove references from the naming service. 

\begin{lstlisting}[frame=tb]{}
${TAO_ROOT}/TAO/utils/nslist/nsdel
\end{lstlisting}
%%$

But \miro also provides a command line option to force the service to
rebind the interface: {\tt -MiroRebindIOR}.

For a complete list of command line option see Chapter
\ref{sec:command-line}.


\section{Starting Client Programs}

\miro comes with a rich set of client programs and utilities that help
you to explore the robots sensory and actuatory devices. Starting a
client program is basically the same as starting a service. Just
change to the directory {\tt \$MIRO\_ROOT/bin} and start whatever
you want. Since these programs aren't bound to specific hardware
devices they can be run on almost any machine.

Note, that to address a service that is registered within other than
the default naming context at the naming service you also have to
provide the {\tt -MiroNamingContext} {\em context} argument at the
command line (in short {\tt -MNC} {\em context}).

%%% Local Variables: 
%%% mode: latex
%%% TeX-master: "miro_manual"
%%% End: 

\chapter{Available Services}
\label{sec:services}

The \miro framework abstracts the robots hardware devices as active
services, that export the sensors and actuators functionality via CORBA
interfaces that can be accessed transparently from other programs,
probably running on totally different machines. Which services there
are available depends on the individual robot type. To provide easy
access to the services interface, each service registers itself under
a standard name at the CORBA Naming Service. To allow multiple
robots to be accessed at a time, each individual robot creates its
one naming context within the Naming Service under which its services
register. The default naming context is {\tt Miro}, but a given robot
should register itself under its own name (for example {\tt stanislav}).
The naming context to use can be specified at the
command line of every service and example program via the
{\tt -MiroNamingContext} option.


\section{Command Line Options}
\label{sec:command-line}

For each service the following command line options can be specified:

\begin{description}
\item[-MiroConfigFile {\em fileName}:]
  The services are widely configurable. These parameters are specified
  in config files within the directory {\tt \$MIRO\_ROOT/etc}. The
  file name is derived from the computers {\tt \$HOST} environment
  variable: {\tt \$HOST}.xml. If you want to specify another
  configuration file, you can do with this command line option.
\item[-MiroNamingContext {\em contextName}:]
  The short form of this option is \textbf{-MNC}.
  As mentioned before all services register their interfaces at the
  naming service within a naming context. The default naming context
  is {\tt Miro}. In multi robot scenarios for instance it would be
  necessary to assign a different naming context to each robot. This
  can be done by this command line option. Note that since it is
  necessary to specify none default naming context names also at the
  client side, this option is also recognized by every client program.
\item[-MiroRebindIOR:]
  To prevent services from different robots to accidentally
  deregister each other from the naming service, an already existing
  entry in the naming service isn't unbound by a service, if it tries
  to register under the same name. Anyhow, if you need to overwrite an
  existing entry, specify this option at the command line.
\end{description}


\section{Individual Interfaces}

In the following, we list the different services and their
interface names in the naming service, all interfaces are
registered within the specified naming context in the naming
service. For multi robot scenarios the canonical naming convention
for the individual robots naming contexts is the robots name.
Don't get confused by the fact that many services are collocated
within a single binary, while for others there exists a dedicated
binary. Note that in general the specialized interfaces of
individual robots register at the Naming Service under their
ancestors name.  The interface documentation can be looked up
within the \miro reference and in the \miro online documentation
as well.  For some services specializations for individual robot
types do exist. Those are described in the subsequent sections.

\begin{description}
\item[Odometry:]
  The motion service registers the {\tt Miro::Odometry} interface as
  {\tt Odometry}. It encapsulates
  the odometry (dead reckoning) sensor.
\item[Motion:]
  The motion service registers the {\tt Miro::Motion} interface as
  {\tt Motion}.
\item[RangeSensor:] The range sensor interface is the general
  abstraction of all range sensor devices, such as sonar, laser range
  finders etc. It offers a method for querying the sensors physical
  configuration as well as the latest sensor reading.  See the {\tt
    QtRangeSensor} program in the utils directory for an example on
  how to use this information.  The range sensor interface becomes
  registered at the naming service under the name of the actual
  sensory device.
  \begin{description}
  \item[Sonar:] The sonar service registers as {\tt Sonar}.  It
    provides an interface to the very common ultrasonic sensors for
    robots.
  \item[Infrared:] It registers as
    {\tt Infrared}.
  \item[Tactile:] Some robots have bumpers as some sort of ''it's too
    late'' sensors.  Still it is better to stop when you crashed than
    just moving on. It is registered as {\tt Tactile}.  Since it can
    be interpreted as a (very limited) range sensor, it is also a {\tt
      Miro::RangeSensor}.
  \item[Laser:] The laser scanner service registers as {\tt Laser}. It
    provides very accurate 180� distance measurements, with up to
    48Hz.
  \end{description}
\item[Stall:]
  Similar to the bumpers, the stall detection monitors the robots
  motion and detects, when the robot is stuck in its movement. The
  {\tt Miro::Stall} interface registers at the naming service as
  {\tt Stall}.
\item[Video:] The video service registers the {\tt Miro::Video}
  interface as {\tt Video}. It provides camera images at a rate of up
  to 25 images per second captured by frame grabber cards. Currently
  supported are {\tt matrox meteor} and cards that are supported
  by {\em video for linux} (bttv8).  Since the bandwidth needed for
  uncompressed image data transportation exceeds the capabilities of
  most todays network devices and even memory copying can introduce
  excessive overhead, the video service enables access
  to the grabbed images via shared memory keys.
\item[PanTilt:] Cameras often are mounted on top of a pan-tilt unit,
  allowing the robot to look sideways while moving in another
  direction. Lightweight versions of such a device are panning or
  tilting units.  The associated services register the
  {\tt Miro::Pan}, {\tt Miro::Tilt} and {\tt Miro::PanTilt}
  interfaces as {\tt Pan, Tilt} and {\tt PanTilt} at the naming
  service.
\item[Buttons:]
  For simple user interaction, some robot types provide push
  buttons. The buttons service registers the {\tt Miro::Buttons}
  interface as {\tt Buttons}.
\item[Speech:]
  To provide a more natural way of communication, some robots are
  equipped with speech synthesizer cards. The speech service registers
  the {\tt Miro::Speech} interface as {\tt Speech}.
\end{description}


\section{Asynchronous Sensory Information}
\label{SEC:ASYNCHRONOUS_SENSORY_INFORMATION}

Another feature of the \miro framework is the asynchronous
distribution of sensory data via the CORBA Notification Service
\cite{OMG:Notify-Spec}. These event channels allow filtered and
priority based event processing for time critical sensory information
distribution under high load.  Therefore in every robots naming
context exists a reference {\tt EventChannel} under which the
robots notification service can be accessed.

The usage of the notification service within the \miro framework is
explained in more detail in \ref{sec:notify}. In the following we give
only a brief description. The notification service allows simple event
filtering on the bases of the domain name and the type name of the
event. This way the individual consumer can easily subscribe for the
messages it is interested in. As the domain name, the naming context
is used within the \miro framework (note that in a multi robot
scenario, an event channel can transport data from multiple robots).
In the following the type names of the events and the data contained
in the event are listed for the individual services:

\begin{description}
\item[Odometry:] The odometry servant and all derived services
  generates periodical {\tt Miro::MotionStatusIDL} events), that
  propagate the robots current position and velocities. The event type
  name is {\tt Odometry}.
\item[RangeSensor:] The range sensor servant and all derived services
  can generate periodical range sensor events. The event type name is
  that of the derived sensor interface (Sonar, Laser etc.). It is also
  returned by the {\tt getScanDescription} method within the {\tt
    ScanDescriptionIDL} struct. The event data is one of three types,
  as also specified within {\tt Miro::SensorDescriptionIDL}.
  \begin{description}
  \item[Miro::RangeScanEventIDL]
    The data of a range sensor that acquires data in a continuous scan
    pushes its data within this struct. (Currently, there does exist
    no actual sensor that behaves in this way, but anyways...)
  \item[Miro::RangeGroupEventIDL]
    The data of a range sensor that acquires data groupwise
    pushes its data within this struct. A range sensor group
    normally is formed by a set of sensors that are mounted on the
    same hight, pointing in different directions. The laser scanners
    are organized that way.
  \item[Miro::RangeBunchEventIDL] The data of a range sensor that
    acquires data in a discontinuous fashion pushes its data within
    this struct. Each sensor reading contains its own group id as
    well as its index within that group. Especially sonar sensor
    are organized that way, since neighbouring sensor mustn't be
    fired simultaneously to minimize crosstalk.
  \end{description}
  The event type name is that of the actual sensor: {\tt Sonar}, {\tt
    Laser}, {\tt Tactile}, etc. as also specified within {\tt
    Miro::SensorDescriptionIDL}.

\item[Sonar:] The sonar service generates range sensor events with the
  type name {\tt Sonar}. Note that the sonar sensors are fired
  interleavingly to avoid crosstalk between them. Therefore the
  payload of an event emitted by a sonar device is normally a
  \texttt{RangeBunchEventIDL}. That way you have to analyze multiple
  {\tt Sonar} messages to get a full sonar scan.
\item[Infrared:] The infrared service generates range sensor events
  with the type name {\tt Infrared}.
\item[Tactile:] The tactile service generates range sensor events with
  the type name {\tt Tactile}. Due to the hopefully low frequency
  at which those events occur it might be better to listen to the
  {\tt Tactile} events, instead of polling the tactile status. Tactile
  devices usually contain a \texttt{RangeBunchEventIDL} struct.
\item[Stall:] The stall service also emits events. These contain the
  type name {\tt Stall}. For the event data see the individual
  robots descriptions.  Due to the hopefully rare occasions at which
  those events occur it might be better to listen to these events,
  instead of polling the stall interface.
\item[Buttons:] They also provide the event triggered communications
  model that uses the event type {\tt Button}. The events payload is
  of type {\tt ButtonStatusIDL}.
\end{description}


\section{\bXXI}

The \bXXI robot has a large number of available sensors and actuators. The
relevant binaries for its specialized services are:

\begin{description}
\item[B21Base:] This is a collection of the robots main services. This
  binary incorporates the motion service as well as the sonar, tactile
  and infrared services. Furthermore it gives access to the four
  colored buttons, that are mounted on top of the robot. The
  corresponding interfaces are:
  \begin{itemize}
  \item {\tt Miro::B21Motion} \\
    The {\tt Miro::B21Motion} interface is derived from the
    {\tt Miro::Motion} interface and gives access to the
    specialized movement commands of the B21 robot. It registers under it
    ancestors name {\tt Motion}.
  \item {\tt Miro::Sonar} \\
    The derived {\tt Miro::RangeSensor} interface registers itself
    as {\tt Sonar}.
  \item {\tt Miro::Infrared} \\
    The derived {\tt Miro::RangeSensor} interface registers itself
    as {\tt Infrared}.
  \item {\tt Miro::Tactile} \\
    The derived {\tt Miro::RangeSensor} interface registers itself
    as {\tt Tactile}.
  \item {\tt Miro::B21Buttons} \\
    The B21 has four buttons, that can be used as a simple user
    interface. The {\tt Miro::B21Buttons} interface is supported by
    this service. It sends {\tt Miro::ButtonStatusIDL} events, if a
    button is pressed or released.
  \end{itemize}

\item[SickLaserService:]
  This is the service to access the SICK laser scanner. It supports
  the {\tt Miro::RangeSensor} interface.
\item[DtlkSpeech:]
  This service lets you control the DoubleTalk Speech cards. The
  interface is named {\tt Miro::Speech}
\item[VideoService:]
  It gives you access to the two frame grabbers of the robot. It
  supports the {\tt Miro::Video} interface.
\item[DpPantilt:]
  It supports the {\tt Miro::DirectedPerceptionPanTilt} interface,
  which is derived from {\tt Miro::PanTilt}. It therefore registers
  itself as {\tt PanTilt}
\end{description}

Note that the B21 has two internal computers, to which the
different sensors and actuators are attached. You have to start
the relevant binaries at the correct computer. The
sensory/actuatory distribution in our laboratory is as follows:
Start B21Base and LaserService on the left computer, VideoService
and DpPantilt and DtlkSpeech on the right.


\section{\sparrow}

The \sparrow soccer robots have a similar diversity of sensors and
actuators, but for efficiency and internal design reasons, they are
grouped together in less binaries:

\begin{description}
\item[SparrowBase:]
  Within this file the following services are grouped together:
  Motion, Stall, Kicker, Buttons, Sonar, Infrared and Pantilt. The
  corresponding interfaces are:
  \begin{itemize}
  \item {\tt Miro::SparrowMotion} \\
    The {\tt Miro::SparrowMotion} interface is derived from the
    {\tt Miro::Motion} interface and gives access to the
    specialized movement commands of the \sparrow robot. It registers
    under it ancestors name {\tt Motion}.
  \item {\tt Miro::Stall} \\
    The {\tt Miro::Stall} interface is registered as {\tt Stall}.
  \item {\tt Miro::Kicker} \\
    Due to its purpose as football robot, the \sparrow robot has a
    kicker, that can be accessed via this interface. It registers at
    the naming service as {\tt Kicker}.
  \item {\tt Miro::Buttons} \\
    The {\tt Miro::Buttons} interface registers as
    {\tt Buttons}.
  \item {\tt Miro::Sonar} \\
    The {\tt Miro::RangeSensor} interface registers itself
    as {\tt Sonar}.
  \item {\tt Miro::Infrared} \\
    The {\tt Miro::RangeSensor} interface registers itself
    as {\tt Infrared}.
  \item {\tt Miro::PanTilt} \\
    The {\tt Miro::SparrowPanTilt} interface is registered as
    {\tt PanTilt}.
  \end{itemize}

  Note that the
  stall service sends {\tt Miro::SparrowStallIDL} events.
\item[VideoService:]
  It gives you access to the frame grabber of the robot. It supports
  the {\tt Miro::Video} interface.
\end{description}


\section{Pioneer}

Due to the limited variety of sensors and actuators on the Pioneer robots, there exists only
the most fundamental services for this robot.

\begin{description}
\item[PioneerBase:]
  The following services are grouped together: Motion, Stall and Sonar.
  The corresponding interfaces are:
  \begin{itemize}
  \item {\tt Miro::PioneerMotion} \\
    The {\tt Miro::PioneerMotion} interface is derived from the
    {\tt Miro::Motion} interface and gives access to the
    specialized movement commands of the Pioneer robot. It registers
    under it ancestors name {\tt Motion}.
  \item {\tt Miro::PioneerStall} \\
    The {\tt Miro::PioneerStall} interface (derived from {\tt
      Miro::Stall}) is registered as {\tt Stall}.
  \item {\tt Miro::Sonar} \\
    The {\tt Miro::RangeSensor} interface registers itself
    as {\tt Sonar}.
  \end{itemize}

  Note that the stall service sends {\tt Miro::PioneerStallIDL} events.
\end{description}


\section{Frame Grabbers}

The {\tt videoService} provides an interface to several frame grabber
cards. All BTTV cards that are supported by the video for linux
project are also supported here (namely the Bt848/849/878/879 based
frame grabbers). The Matrox Meteor frame grabbers cards are also
supported, but note that kernel drivers are only available for the
older 2.2.x kernel series.

The idl-interface that is used by the {\tt videoService} is {\tt
  Miro::Video}. The Server registers itself as {\tt Video} in the
Naming Service. You can get the last grabbed image immediately by use
of the {\tt getImage()} method or wait for the next available image
with the {\tt getWaitImage()}. What you get is a handle to the image
that simply points to a memory range, where the color values for the
image (or the intensity value for each pixel if it is a grayscale
image) are listed for each single pixel. Additionally you get the
image dimensions, an associated time stamp and the robots position at
the moment of recording the picture.

As always there is a configuration section in the xml-file called {\tt
  video}, where the following parameters can be used:

\begin{description}
\item[device:] This section describes, which device-name the frame
  grabber card uses.
\item[grabber:] Determine which card your robot uses, either a {\tt
    bttv} based card or a matrox {\tt meteor} card.
\item[format:] With the format parameter you can choose, which
  standard you camera uses. Either the (european) {\tt pal} format,
  the (american) {\tt ntsc} format or the (asiatic) {\tt secam}
  format. With the {\tt auto} option, the format is choosen
  automatically.
\item[source:] This parameter defines, from which input source the
  server should read. Possible choises are: {\tt composite1}, {\tt
    composite2}, {\tt composite3}, {\tt svideo} and {\tt tuner}.
  %%{\tt usb} or {\tt 1394}.
\item[palette:] Palette can be one of the following: {\tt gray}, {\tt
    rgb}, {\tt rgba}, {\tt bgr} or {\tt bgra}. \\
  {\tt rgb} or {\tt bgr} indicate the 24-Bit mode, whereas {\tt rgba} and
  {\tt bgra} are 32-Bit modes.
\item[subfield:] With this parameter you can choose which subfields to
  use: {\tt odd}, {\tt even} or {\tt all}.
\item[width:] The {\tt width} that the image should have.
\item[height:] The {\tt height} that the image should have.
\item[connections:]
\item[buffers:]
\item[byteorder:] {\tt 0}
means native mode, {\tt 1} means rgb.
\item[pixelsize:] This
parameter specifies the size in bytes of each pixel (1, 3 or 4).
\item[upsidedown:] On some systems the camera is mounted upside
down. Set this parameter to true if this is the case.
\end{description}

Note, that not each card supports each option and that not every
combination of options can be used for each card. If in doubt,
refer to the technical manual of your frame grabber card.

The following example shows the configuration used for our hardware:

\begin{lstlisting}[frame=tb]{}
<video>
  <device>/dev/v4l/video0</device>
  <grabber>bttv</grabber>
  <format>pal</format>
  <source>composite1</source>
  <palette>bgr</palette>
  <subfield>odd</subfield>
  <width>384</width>
  <height>288</height>
  <connections>16</connections>
  <buffers>16</buffers>
  <byteorder>1</byteorder>
  <pixelsize>3</pixelsize>
</video>
\end{lstlisting}


\section{SICK Laser Scanner}

The {\tt sickLaserService} provides an interface to a laser range
finder of type Sick PLS.
%%pls200 (TODO: check this?).
This sensor is connected to the controlling PC via a serial line. It
is delivered either with a RS-232 or a RS-422 compliant interface.

The idl-interface implemented by the {\tt sickLaserService} is
{\tt Miro::Laser}, which is derived from
{\tt Miro::RangeSensor}. The server can be driven in two different
modes, either the sensor automatically provides a scan about 40 times
a second, or every scan is polled, which results in a maximum scan
frequency of 10 times a second. The second mode also allows to choose
a different (i.e. lower) scan frequency. At the moment no other
special features of the Sick PLS200 are supported.

The configuration of this service is provided in xml, within the
section {\tt sick}. The following parameters are understood:

\begin{description}
\item[device:]
  The filename of the device which should be used.
\item[baudrate:]
  The baudrate for communication with the sensor,
  possible values are: 9600, 19200, 38400 and 500000.
\item[stdcrystal:]
  For achieving the unusual baud rate of 500000
  baud, we use a serial adapter, that is equipped with a 16MHz
  crystal, instead of the standard 14.??MHz. Due to this reason, the
  standard baud rates must be generated by modifying the divisor of
  the UART. If you have an off-the-shelf serial adapter, use
  {\tt true} here, you will probably not be able to use 500000 baud
  then.
\item[continousmode:]
  If {\tt true} is provided for this
  parameter the sensor will provide data automatically, about 40 times
  a second. If {\tt false} is used, the service will poll a
  measurement after an interval that can be specified with the
  following parameter.
\item[pollintervall:]
  The interval between two consecutive
  measurements, measured in microseconds. This parameter has no effect
  if continousmode is set to {\tt true}.
\end{description}

The following example shows the configuration used for our hardware:

\begin{lstlisting}[frame=tb]{}
<sick>
  <device>/dev/laser</device>
  <baudrate>500000</baudrate>
  <continousmode>true</continousmode>
  <stdcrystal>false</stdcrystal>
  <pollintervall>100000</pollintervall><!-- pollinterval in usec -->
</sick>
\end{lstlisting}

As the laser provides a specialization of a range sensor, we also
provide a scandescription for it. See section \ref{scandescrption} for
details.


\subsection{Possible Problems}

Due to the unusual and high data rate of 500000~baud a special serial
interface card is required. The card recommended by Sick provides a
16550A compatible interface, for that reason the current
implementation is able to use the card through the standard linux
serial device drivers, which keeps the implementation of the service
more portable. Unfortunately this introduces a big drawback: receiving
up to 50000 characters per second, results in at least $50000/16=3125$
interrupts per second, this is an enormous load. It is essential, that
every interrupt is handled on time, because the sensor does neither
support hardware, nor software flow control, and provides a packet
with data every 25~ms. If you have the chance to use an interface card
with a larger FIFO buffer than the standard 16 bytes, this will
provide a large improvement in performance, stability and system
reactivity.  We solved most of our stability problems, that were
caused by lost packages without a larger FIFO: We enabled irqs
throughout IDE interrupt processing (this may be extremely dangerous,
depending on your configuration, see man page of {\em hdparm}). Another
solution could be using {\em irqtune} to give the serial irq a higher
priority, or using {\em RT-Linux}.

\textbf{Due to the adaptations of the divisor for the serial line,
the {\tt SickLaserService} has to be suid root. If you use a
standard crystal, and you do not use 500000~baud, this is not
necessary.}


\section{DoubleTalk Speech Card}


\section{Directed Perception Pantilt}

%\section{Higher-level Servers}

%In the current version of \miro, no high-level services are included.

%\subsection{MapServer}

%\subsection{RegionServer}

%\subsection{SaliencyServer}

%\subsection{ClassificationServer}


%%% Local Variables:
%%% mode: latex
%%% TeX-master: "miro_manual"
%%% End:

\chapter{Test and Example Programs}
\label{sec:clients}

The \miro repository comes with different kinds of test and
example programs that are scattered over four different directories:

\begin{description}
\item[examples]
  These programs show the easiest way of using a
  specific service. In most cases it's just getting data and printing
  it to {\tt cout}.
\item[utils] These are programs designed to fasciliate the handling of
  your mobile robots. For example, programs located in this directory
  provide a graphical view of the services output or state.
\item[tests]
  These programs are designed to reliably
  test an entire interface. They often come with a simple character
  dialog and the possible selections rely directly to the interfaces
  as described in the auto-generated online help.
\item[performance-tests]
  They measure the performance of individual
  sensory or actuatory devices. Mostly the throughput.
\end{description}


%\section{Simple Examples}

%\subsection{baseClient}
%- baseServer must run

%\subsection{sonarClient}
%- baseServer must run

%\subsection{irClient}
%- baseServer must run

%\subsection{laserClient}
%- baseServer and laserServer must run


%\section{Graphical Sensor Clients}

%\subsection{qtSonar}

%- screenshot

%- baseServer must run

%\subsection{qtLaser}

%- screenshot

%- baseServer and laserServer must run



%\section{Test Suite Programs}



%%% Local Variables: 
%%% mode: latex
%%% TeX-master: "miro_manual"
%%% End: 

\chapter{Writing Your First Programs}

This chapter tries to help you with the first steps of writing
programs that use the \miro framework. The CORBA environment and the
\miro framework seem to raise the bar for an easy entry into robot
programming. While this can hardly be denied they facilitate
tremendously the task of writing distributed programs. And since robot
control software is inherently distributed (ever thought of multiple
robots?) it seems the only way to go.

As you will see, most of the distributed programming complexity is
initially hidden from the user:

\begin{itemize}
\item The programmer simply calls methods of the devices interface.
\item The programmer communicates via proxy object to the remote
  service. The most tricky part is getting the object reference.
\item The distributed environment is transparent. The
  remote method invocation (RMI) is hidden entirely.
\end{itemize}


\section{Makefiles}
\label{sec:makefile}
In contrary to earlier versions of \miro, it is no longer based on the
ACE/TAO Makefile hierarchy, but uses the GNU autotools for a simpler
and more portable way to generate the Makefiles (beside that, this
inherently adds new nifty features, like e.g.\ the automatically
distribution via a simple \texttt{make dist}). Nevertheless it is
still possible to use the ACE/TAO style of writing Makefiles in
derived projects, as we didn't delete the \miro include rules and
macros (but please note that we do not maintain them any longer and
they may be deleted in future version).

Makefile templates reside in the directory {\tt
  \$MIRO\_ROOT/templates}, that should facilitate the creation of
further subprojects within the \miro directory structure but can just
as well be used to start new projects outside the {\tt \$MIRO\_ROOT}
directory. The makefiles are designed to either build subdirectories,
libraries or binaries. To use a template, simply copy it into the
corresponding directory under the name {\tt Makefile.am} and adapt it
as explained in the following sections. Additionally, you have to put
a line in \texttt{configure.ac} to the macro named
\texttt{AC\_CONFIG\_FILES} to let automake and autoconf know, that they
have to build a \texttt{Makefile} (\texttt{Makefile.in} respectively).

\subsection{Makefile.am.dir}
This Makefile starts subbuilds in the specified directories. It is an
example for conditional compilation (depending on the
configuration). If you don't have conditional branches, you can left
the second variable empty or reduce the Makefile.am to a single line,
like e.g. \texttt{SUBDIRS = firstDir secondDir}.

\subsection{Makefile.am.bin}
This Makefile builds a single or several binaries. The name of the
binaries to be built are specified in the variable {\tt bin\_PROGRAMS}.
There have to be an extra line for each binary to be build starting
with the name of the binary itself and an additional
\texttt{\_SOURCES}-suffix. The sources are specified with their suffix
(.cpp). Please also list all the header files here, otherwise they
don't get distributed by an call to \texttt{make dist}. If the
programs need to be linked against several libraries, there is again
one line per binary containing the name of the binary plus an
additional \texttt{\_LDADD}-suffix.

\subsection{Makefile.am.qt}
This is a special Makefile.am template for Qt programs. For
simplicity, the Makefile is designed for a single binary per
directory, but beside that, most of the above section also applies
here. The only exception is, that we have to distinguish between
normal source files (variable \texttt{sources}) and files (variable
\texttt{tomocsources}) that need to be compiled with Qt's Meta Object
Compile (\texttt{moc}).

\subsection{Makefile.am.lib.a}
This Makefile builds a static library. The name of the library can be
specified in the variable \texttt{lib\_LIBRARIES}, the files building
the library in the variable \texttt{sources} (with the .cpp suffix; or
.xml in case of an xml-file for \miro's parameter framework). Headers
that are not automatically considered because they don't have a
corresponding cpp-file should be specified in the variable
\texttt{extraheader}. Don't forget to name the prefix for the
\texttt{\_SOURCES}-macro with the name of the resulting library.
Therefore please substitued all occuring dots with underscores (e.g.\ 
the usual \texttt{.a} a the end of an static library).  Additionally,
replace the prefix for the macro \texttt{\_include\_HEADERS} with the
name of the directory, to which the header files should be installed
(see also the description of \texttt{Make-rules} below).

\subsection{Makefile.am.lib.so}
Makefiles for shared libraries look quite the same as for static
libraries, to most of the above applies here too. The only difference
from a users perspective is the name of the variable
\texttt{lib\_LTLIBRARIES} where you name your library (note the
\texttt{LT} within the variable (and the targets at the end of the
\texttt{Makefile.am}) which indicates that the library is build using
the libtool package).

If you programs or libraries need to be compiled differently depending
on configuration options, please have a look at an existing
\texttt{Makefile.am}, e.g.\ in
\texttt{\$MIRO\_ROOT/\-examples/\-behaviours/\-engine/} or
\texttt{\$MIRO\_ROOT/\-src/\-miro/}.

\subsection{Make-rules}
Last you may have a look at the \texttt{Make-rules} file within the
\texttt{\$MIRO\_ROOT/config/} directory. This file is included into each
\texttt{Makefile.am}, so rules or variables that are needed globally,
can be inserted here. If you want to define additional installation
subdirectories for the header files, you can see examples here.


\section{A Simple Sonar Client}

To discuss things using actual code, let's look at the simple task of
obtaining data from a sonar sensor device. Since this is a range
sensor, the sonar is queried via the generalized range sensor
interface. The only difference in querying an infrared or a sonar
device is the name under which these sensors are registered within the
naming service.

\lstinputlisting[frame=tb, caption={examples/sonar/SonarPoll1.cpp}]{SonarPoll1.cpp}
\label{lst:SonarPoll1}

A step by step walk trough the code:

\lstinputlisting[frame=tb, first=1, last=1]{SonarPoll1.cpp}

The file \lstinline!miro/Client.h! contains the definition of the class
\lstinline!Miro::Client!. See below.

\lstinputlisting[frame=tb, first=2, last=2]{SonarPoll1.cpp}

The file \lstinline!miro/RangeSensorC.h! provides the classes for the interface of the
sonar service. The 'C' at the end of the name stands for client. This
file is automatically generated from the idl-description of the interface.

\lstinputlisting[frame=tb, first=6, last=6]{SonarPoll1.cpp}

A standard main function is used.

\lstinputlisting[frame=tb, first=8, last=8]{SonarPoll1.cpp}

This class just wraps the few lines that are necessary at the beginning of any CORBA
application. It will therefore be normally instantiated in all \miro
applications first. This is a simple helper class that sets up the
CORBA environment in a standard way. Nothing tricky there.  It
initializes the orb etc. The call uses the command line arguments for
finding the naming service, and parsing other commands...

\lstinputlisting[frame=tb, first=9, last=10]{SonarPoll1.cpp}

This initializes a proxy object to the sonar service. This proxy
object is used as if it was the sonar service itself. Let's look
at the call a bit closer. \lstinline!Miro::RangeSensor_var! is the
CORBA equivalent of a standard C++ auto pointer for a
\lstinline!Miro::RangeSensor! object. That is the proxy becomes
automatically destroyed if the pointer goes out of scope. The
method {\tt resolveName()} is a template member function. Its
argument specifies the name that shall be resolved in the default
naming context of the naming service (see Chapter
\ref{sec:using}). The \lstinline!Miro::RangeSensor! in French
brackets specifies the type of reference that shall be returned by
the call. Note that this will only succeed if the reference stored
under the name of Sonar refers indeed to a {\tt Miro::RangeSensor}
object or to a derived ancestor of this class.

\lstinputlisting[frame=tb, first=11, last=11]{SonarPoll1.cpp}

The RangeSensor interface returns a pointer to the sensor scan, which
the caller obtains ownership of. So we use another auto pointer to
hold the return value of the call and to prevent us from memory leaks.
The IDL in the type name reflects the fact, that this is a IDL defined
data type. (This is just a naming convention of \miro that shall help
you to trace the roots to the documentation.) In the CORBA mapping for
C++ these types are mapped to plain C structs: No methods, no
inheritance, just public data members.

\lstinputlisting[frame=tb, first=13, last=13]{SonarPoll1.cpp}

This gets a sonar scan from the service. Note, that we do not see where
the service runs.

\lstinputlisting[frame=tb, first=15, last=20]{SonarPoll1.cpp}

These lines are used to print the received sonar scan to the
console.

%\lstinputlisting[frame=tb, first=22, last=22]{SonarPoll1.cpp}

%Ends the program. Returning zero indicates, that no error occurred.

%%\bigskip
%%- Note the {\tt Miro::} prefixes which define the namespace ...


\section{Using Namespaces}

By inserting the lines

\lstinputlisting[frame=tb, first=6, last=9]{SonarPoll2.cpp}

we map the relevant ``Miro'' data types from the \lstinline!Miro!
namespace into our global namespace. This shortens the type name
specifiers, but also hides from where they are coming from. Note that
a simple

\begin{lstlisting}[frame=tb]{}
using namespace Miro;
\end{lstlisting}

would do the same trick. But mapping a namespace completely is
generally not a good idea, since it tends to produce name conflicts
and secondly code reviewers can trace the origin of the types less
easily. Be especially careful in header files. You are flattening the
namespace for everyone that has to include your header file, which can
lead to bad surprises.

See the following listing which shows the same code again, but without
the \lstinline!Miro::! prefixes. The initial lines containing the includes
are skipped.

\lstinputlisting[frame=tb, first=6, caption={examples/sonar/SonarPoll2.cpp}]{SonarPoll2.cpp}
\label{lst:SonarPoll2}


\section{Adding Exception Handling}

Handling of error conditions itself is error prone. Exceptions make
the handling of error conditions easier, but errors are errors and
therefore stay somehow nasty. What's quite easy to achieve by the use
of exceptions is to print some diagnostic output and exit instead of
gracefully segfaulting. This is done in this example by enclosing the
code in the main function in a try/catch block.

Since \miro defines ostream operators for every \lstinline!Miro::...IDL!
type in \lstinline! miro/IO.h!, we include that file and abandon the
handcrafted streaming of the sonar data.

Note the different kinds off exceptions. There are \miro
exceptions. Exception types defined in \miro are beginning with a
big E as a naming convention. These indicate problems on the
service side, like hardware problems (may be the batteries?), bad
service calls (trying to accelerate the robot to warp 1?) or load
problems. Then there are CORBA exceptions. Those occur if there
arise some communication problems: A service went down, the robot
is loosing the radio ethernet connection etc. Since all \miro
exceptions derived from \lstinline!CORBA::UserException!, those
also are catched within the first \lstinline!catch! block of the
example code. You do not have to catch all exceptions. An uncaught
exception will lead to program termination. Just as if you'd catch
them at the end of main and exit...

Note also, that the instantiation of Client isn't within the try/catch
block. This is intentionally. Exceptions that can arise in the
construction of a Client instance are CORBA exceptions. The ostream
operators for CORBA exceptions used by TAO need an ORB
instance. Since the ORB is instantiated within the Client class, it
will not exist after destruction of the Client --- and this would be
done at the end of the try block. Therefore we catch exceptions within
the constructor of Client, print them to stderr and exit. There is
little to do for the client program anyhow, if it can't access the
services.

\lstinputlisting[frame=tb, first=14, caption={examples/sonar/SonarPoll3.cpp}]{SonarPoll3.cpp}
\label{lst:SonarPoll3}

\section{An Asynchronous Sonar Client}
\label{sec:notify}

By now we were actively requesting for the data of a service. But
think of polling for tactile events that way. Having a good
collision avoidance, those events should hardly ever occur.
Nevertheless, as soon as there is a tactile signalling a
collision, the robot should immediately react to this event. By
polling it would have to call the tactile interface thousands of
times, just not to miss the one event it can't effort to miss. And
since the sensors are actively collecting their data, shouldn't
they be able to trigger the data processing within the robot? - Oh
yes, they can.

For this purpose the notification framework within \miro does
exist. It is based on the CORBA Notification Service
\cite{OMG:00-5}, and precustomized by some utility classes. They
enable clients to subscribe to arbitrary events of a notification
channel. The data gets pulled to them as soon as it becomes
available at the producer (e.g. a range sensor device).

To show how asynchronous event processing works within \miro, lets
look at a small example. First we look at the code to handle the
events.

\lstinputlisting[frame=tb, first=24, last=53, caption={examples/sonar/SonarNotify.cpp}]{SonarNotify.cpp}
\label{lst:SonarNotify}

A step by step walk trough the code:

\lstinputlisting[frame=tb, first=24, last=24]{SonarNotify.cpp}

The class \lstinline!SonarNotify! provides a callback for the event
channel. It is derived from \lstinline!Miro::StructuredPushConsumer!,
which handles the registration at the notification service etc. All
the initialization and registering is done within the constructor of
this super class. The method \lstinline!push_structured_event! is the
callback that is called from the notification service.

\lstinputlisting[frame=tb, first=31, last=31]{SonarNotify.cpp}

The \lstinline!EventTypeSeq! class is for specifying the events that you
want to be subscribed for. It is an incremental protocol. So you make
a list of events you want to receive from now on and another with the
events you wish to no longer become informed about. Since we are just
subscribing, the second list is empty.

\lstinputlisting[frame=tb, first=33, last=35]{SonarNotify.cpp}

Events are subscribed by \lstinline!domain_name! and
\lstinline!type_name!. Never forget to set the length of a
sequence explicitly. Specifying the length as a constructor
parameter just reserves the number of elements, the length of the
sequence is still zero.

\lstinputlisting[frame=tb, first=37, last=37]{SonarNotify.cpp}

Tell the consumer admin what we want to subscribe for.

\lstinputlisting[frame=tb, first=48, last=49]{SonarNotify.cpp}

The payload of a structured event is contained within the
\lstinline!remainder_of_body! field with is of type
\lstinline!CORBA::Any!.  Therefore you can get a const pointer to the
data with overloaded the \lstinline!operator >>= ()!. The return value
is a \lstinline!CORBA::Boolean! indicating success of the operation.
Even though a \lstinline!CORBA::Any! can hold any IDL defined struct,
you can only extract type \lstinline!T! from a \lstinline!CORBA::Any!
if it actually contains an object of type \lstinline!T!.

\lstinputlisting[frame=tb, first=54, caption={examples/sonar/SonarNotify.cpp}]{SonarNotify.cpp}
\label{lst:SonarNotify2}

Lets now have a look at the \lstinline!main! function:

\lstinputlisting[frame=tb, first=57, last=57]{SonarNotify.cpp}

Since a push consumer is called (pushed) by the event producer, it
is actually a server instead of a client. Therefor the
\lstinline!Miro::Server!  class is instanciated. It performs the
necessary calls to init the CORBA environment. It is basically the
same as the \lstinline!Miro::Client!  but we also need a POA to
register the consumer object.

\lstinputlisting[frame=tb, first=59, last=59]{SonarNotify.cpp}

Just like the former \lstinline!Miro::RangeSensor! interface, we
resolve the \lstinline!EventChannel! by name at the naming service.

\lstinputlisting[frame=tb, first=61, last=61]{SonarNotify.cpp}

Instanciate the push consumer. It does all the necessary
initialization within its constructor.

\lstinputlisting[frame=tb, first=63, last=63]{SonarNotify.cpp}

Now we enter the CORBA event loop. This will not return, until the
process is signaled by SIGINT or SIGTERM.

%%% Local Variables:
%%% mode: latex
%%% TeX-master: "miro_manual"
%%% End:

\chapter{Parameter Sets}
\label{SEC:PARAMETERS}

                                % problem

Robot control programs and robotic algorithms are full of parameters.
Different robot types vary in shape and sensor configuration, but also
different robots of the same type tend to vary slightly. For instance
some of them are equipped with further sensory or actuatory devices
not present in the standard configuration. Also damage and repair of
parts of a robot during its lifetime result in an individual robot,
which has its own unique configuration.

                                % examples

Furthermore the environment provides a complete set of
parameters. Lighting conditions, the shape and location of rooms and
corridors, the position of obstacles like tables or stairs. The list
could be continued indefinitely. These parameters vary between
different environment, but also tend to change slightly over
time. Light bulbs are changed, cupboards added, tables moved etc.

A third source of parameters are defined by the task, the robot has to
perform or the scenario, the robot is designed to operate in.
Generalized tool boxes like planners or knowledge bases have to be
populated with the actions and objects of relevance and a reactive
execution engine need to be configured with what actions to take in
which situation.

                                % summary

The configuration of these parameter sets can be done either by hand
or tool supported or fully autonomously. Nevertheless it is a crucial
part of the setup of robot for every scenario. The handling of those
parameter sets therefore has to be well defined in order to keep the
robots control software maintainable and adaptable to new tasks.

                                % solution aproach

In this chapter we will therfore introduce the parameter framework of
\miro.                                % organization of the chapter
The remaining of this chapter is organized as follows.  First the
general concept of parameter sets and configuration files is
sketched. Then the
location and naming of the robot configuration files of \miro and how to
create a configuration for an individual robot will be explained. In
the following section the description and implementation of
parameter sets and the mapping to the actual parameter files will
be discussed. In section \ref{SEC:DESCRIPTION_SYNTAX} the syntax and
semantics of the description language will be explained.

\section{Overview}

Parameter sets are organized as follows.
 
From a user perspective parameter sets are described by an XML file,
that the service parses on startup, to initialize its parameter set
classes. Parameters in general are strictly typed name value pairs.
They are grouped into sections within the file. Also, they can be
nested, structured types. The parameter names start with an upper case
initial letter.

From a programmers perspective, a parameter set is represented by a
class within C++ with all member variables declared public, just like
a C struct. The member variable names begin with a lower case initial
letter. Names consisting of multiple words are concatenated, the
beginning of a new word is indicated by an upper case letter.
Parameter classes support single, public inheritance. The following
methods are declared:
\begin{description}
\item[Default Constructor]
  Initializes all member variables to their default values.
\item[void operator $\ll =$ (QDomNode)]
  Parses an XML node and sets the member variables accordingly.
\item[QDomElement operator $\gg =$ (QDomElement)]
  Appends the current
  values of the member variables as a child to the XML node.
\item[void printToStream(ostream\&)]
  Prints the current values of the member variables to the
  specified output stream.
\end{description}
Optionally parameter set classes contain a static member method
{\tt instance()}, that returns a pointer to the global heap allocated
instance of the parameter set class.

As parameter set classes have a very generic structure, their coding,
especially the XML and ostream operators can be automated. Therefore
parameter sets normally are described not within a programming
language, but within an XML syntax. The mapping to C++ is then
performed by a small compiler, which emits the necessary code
fragments. The parameter set description can also used by tools for
comfortable editing of configurations by the user. This is done by the
policy editor (see section \ref{SEC:POLICY_EDITOR}) for editing
behaviour parameter sets.

Within \miro this framework is used for two parts of the above
sketched three sources of parameter sets. For the description of the
robots configuration and for the robots task
description. Environmental modeling for different scenarios lies out
of the scope of the problem set addressed by \miro.

\section{Configuration Files}

Every service of \miro tries to find a configuration file on startup
to adapt itself to the individual robots configuration if necessary.
If no configuration file is found, the service has to use its default
settings. The locations the configuration files are expected are:
\begin{enumerate}
\item {\tt \$MIRO\_ROOT/etc}
\item The current directory.
\end{enumerate}
The name of the configuration file is expected to be the host name of
the computer the service is running on. The file extension is .xml:

{\tt \$HOST.xml}

For example in the case of our famous Sparrow99 soccer robot goal
keeper named haeltnix, the \miro service SparrowBase, started on the
computer mounted on the robot will try to load the configuration file:

{\tt \$MIRO\_ROOT/etc/haeltnix.xml}

The default name and location of the configuration file can be
overwritten by the command line option {\tt -MiroConfigFile} {\em
  filename} or {\tt -MCF} for short. The search order of directories
stays the same.  Additionally there can also be an absolute path
specified.

In {\tt \$MIRO\_ROOT/etc} there are sample configuration files for
every robot model currently supported by \miro. They are named
{\em model}.{\tt xml}. Copy the respective configuration file of 
your robot model to {\tt \$HOST.xml} to get a configuration file to 
start with. XML files are text files and quite easily
readable. Therefore opening the configuration file into your favorite
editor and changing the {\em value} attributes to different settings
should be sufficient to solve the first few configuration problems
like differing device names etc. For anything else, read on.

\subsection{Configuration File Syntax}

The available tags and there nesting capabilities look like follows:

$<$!--MiroConfigurationDocument--$>$

\begin{description}
\item[configuration]
  One file always describes one entire configuration. A
  configuration describes all the services parameters, that can be
  executed on the particular robot.
  
  There can only be one configuration tag per file.
  
  \begin{description}
  \item[section] Each component of \miro can have its own section
    within a configuration file. Configuration file sections are kind
    of a weak concept. For instance the vision system is considered to
    be a component within a \miro configuration file. Normally each
    service defines its own section within the configuration file.
    
    There can be multiple section tags within a configuration.
    
    Attributes:
    \begin{itemize}
    \item \textit{name} The name of the section.
    \end{itemize}
    
    \begin{description}
    \item[parameter] Parameters of a component are grouped and
      described as set of parameters that belong together. Note that
      within the configuration file \miro takes a hardware centric
      view. That is, sensors and actuators that result in multiple
      service interfaces, are most likely to be grouped into one
      component if they are controlled by the same low level
      controller board. For instance the parameters for the pioneer
      board, that controls the motors as well as the sonars etc., are
      grouped within just one component.
      
      There can be multiple parameter tags within one section.
      
      Attributes: 
      \begin{itemize}
      \item \textit{name} Name of the parameter. On top level, the
        parameter name attribute of the parameter tag denotes the name
        of the parameter type. The parameter type definitions are
        described in section \ref{SEC:PARAMETER_TYPES}.
      \end{itemize}  
      
      \begin{description}
      \item[parameter] Parameters can be specified within the configuration
        file for each component.
        
        There can be multiple nested parameter tags.
        
        Attributes: 
        \begin{itemize}
        \item \textit{name} Name of the parameter.  
        \item \textit{value} The value of the parameter.
        \end{itemize}
        
        In case the parameter describes a structured type instead of a
        basic type (like bool, int, string etc.), parameter tags can
        be nested. In this case the value attribute is ignored.
        
      \end{description}  
    \item[instance] As the name of a top level parameter tag denotes
      the type name of the parameter set, multiple instances of
      parameters sets in one section become ambiguous. So this tag
      solves the problem. It has two mandatory attributes:
      
      Attributes: 
      \begin{itemize}
      \item \textit{name} Name of the parameter. In this case, an arbitrary name.
      \item \textit{type} The type name of the parameter.
      \end{itemize}
      
      Appart from this, instance tags behave identical, to toplevel
      parameter tags. That is, the values of the parameter set
      elements are described in nested parameter tags.
    \end{description}
  \end{description}
\end{description}

\subsection{Example Configuration}

See {\tt \$MIRO\_ROOT/examples/params/TestConfig.xml}
for a small initial example.

TODO: include and comment the example code here.

Section tags are structurising elements and component parameter tags
are mainly defined by the underlying hardware devices. But what
defines which parameters are accepted within by which component and of
what types they are? While the configuration files contain comments to
explain the purpose of the different parameters, there do also exist
formal definitions of the components parameter sets, as described
within the next section.

\section{Parameter Set Generation}
\label{SEC:PARAMETER_TYPES}

Parameter sets, described as component parameters within a configuration file
correspond to strictly typed classes within C++. Those structures can
be initialized by parsing the configuration file. To facilitate the
generation of parameter sets, parameter classes needn't to be
hand coded, but are auto-generated from parameter description files by a
little parameter class compiler ({\tt MakeParams}). The parameter
description files are once again based on the XML syntax.

Furthermore the description files are not only used for the
configuration of the services of \miro but also by other tools within
the \miro framework. The behaviour framework (sec.\ 
\ref{SEC:BEHAVIOURS}) uses this functionality for parameterization of
behaviours and arbiters and the {\tt PolicyEditor} (sec.\ 
\ref{SEC:POLICY_EDITOR} as well as the {\tt ConfigEditor} (sec.
\ref{SEC:CONFIG_EDITOR}) use the parameter description files for GUI
based editing of the behaviour parameters). To facilitate the reuse of
the parameter descriptions by multiple tools, the syntax of the file
is somehow obfuscated by additional tags and attributes.

\subsection{Description File Syntax}
\label{SEC:DESCRIPTION_SYNTAX}

Parameter classes support single inheritance.

$<$!--MiroParametersConfigDocument--$>$

\begin{description}
\item[config]
  A parameters config file describes the types, names and defaults
  for parameter classes as used within the \miro framework. 

  \begin{description}
  \item[config\_group]
    Configurations can be grouped together for various reasons. 
    A group within a description file in general corresponds to
    a section within the configuration file.

    Attributes:
    \begin{itemize}
    \item \textit{name} The name of the config group. This corresponds
      normally to the section used later within the configuration file,
      but for base classes, the config\_group can be chosen
      arbitrarily.
    \end{itemize}

    \begin{description}
    \item[config\_item] A configurable item to built a parameters
      class for. Each config item will be mapped to a parameter class
      by the compiler.

      Attributes:
      \begin{itemize}
      \item \textit{name} The name of the parameter class. The suffix
        Parameters will be added to the name.
      \item \textit{parent} The name of the super class. The suffix
        Parameters will be added to the name.
      \item \textit{instance} Whether or not a singleton instance
        of the parameter class shall be created (sec.\
        \cite{SINGLETON}). Default is {\em false}.
      \item \textit{final} Corresponds to the finalize keyword of
        java. It is only relevant for the editing tools, not for code
        generation. It indicates, that a toplevel instance of this
        parameter can be created.Default is {\em true}.
      \item \textit{dummy} Sometimes you have to fake a parameter
        class to the editing tools that does not exist as such. No
        code will be generated for such a class. Default is {\em
          false}.

      \end{itemize}
      
      \begin{description}
      \item[config\_parameter] A parameter for the class.

        Attributes:
        \begin{itemize}
        \item \textit{name} Guess what.
        \item \textit{type} Allowed basic types are: bool, char, int,
          unsigned int, short, unsigned short, long, unsigned long,
          double and std::string. There are two more
          predefined types: Angle and angle. {\em Angle} corresponds
          to a Miro::Angle instance. It is an angle $\phi$ normalized
          to $-\pi \le \phi \le \pi$ while {\em angle} corresponds to
          a none normalized angle represented by a simple double. Note
          that while the internal representation of angles is in
          radiant, the values specified in the configuration files are
          in degrees. But you can also specify parameter
          classes described within a parameter description file as
          parameter type. Additionally, {\em std::vector$<>$} and {\em
            std::set$<>$} are supported. Note
          however that we are currently lacking support for specifying
          default values for these none basic types.
        \item \textit{default} default value as set in the
          constructor. (optional)
        \item \textit{measure} The measure the type represents, such
          as: mm, msec, �, mm/s, �/s. (optional)
        \item \textit{inherited true/false}
          This attribute is used to 
          overwrite defaults for inherited variables within the
          constructor of the class. If set to true, no new member
          variable will be added to the class, and only the
          attributes name and default will be evaluated.
        \end{itemize}
      \end{description}
    \end{description}
  \item[config\_global]
    Global configuration items those are mainly used for code
    generation. They specify things, that the compiler can not easily
    derive autonomously, such as additional include directives, etc.

    Attributes:
    \begin{itemize}
    \item \textit{name} The name of the config item. There are
      currently three items supported:
      \begin{description}
      \item[namespace] Specifies the namespace of the generated
        code. While multiple parameter classes can be described
        within one file, they all have to reside within the same
        namespace.
      \item[include] Specifies a user include directive.
      \item[Include] Specifies a system include directive.
      \item \textit{value} The name of the namespace, the local or
        global include directive.
      \end{description}
    \end{itemize}
  \end{description}
\end{description}

\subsection{Example Description}

See {\tt \$MIRO\_ROOT/examples/params/Parameters.xml}
for a small initial example.

TODO: include and comment the example code here.

\subsection{MakeParams}

The parameter description compiler has the following command line
parameters.
\begin{description}
  \item[-f $<$file$>$] name of the input file (default is {\em Parameters.xml})
  \item[-n $<$name$>$] base name of the output file (default is {\em Parameters})
  \item[-s $<$extension$>$] extension of the generated source file (default is {\em cpp})
  \item[-h $<$extension$>$]extension of the generated header file (default is {\em h})
  \item[-v] verbose mode
  \item[-?] help: emit this text and stop
\end{description}

\subsection{Example Header File}

See {\tt \$MIRO\_ROOT/examples/params/Parameters.h}
for a small initial example.

TODO: include and comment the example code here.

\subsection{Makefile magic}

To enable automatic generation of the C++ source and header files from
XML parameter descriptions, there does allready exist a rule within
the make files of \miro. That is, for every file named *Parameters.cpp
| h, make checks the existance of the corresponding *Parameters.xml
file and if necessary, invokes MakeParams accordingly. To prevent make
from removing the generated files at the end of the task, you have to
define an additional .PRECIOUS target like this:

.PRECIOUS: MyParameters.cpp MyParameters.h

\section{Config File Editor}
\ref{SEC:CONFIG_EDITOR}

For advanced config file editing, like filter trees of the video image
processing framework \ref{SEC:VIDEO_IMAGE_PROCESSING}, a simple text
editor is a fairly errorprone tool. XML editors do a better job in
this case, but for real comfort, \miro provides the config file
editor.

%%% Local Variables: 
%%% mode: latex
%%% TeX-master: "miro_manual"
%%% End: 

\chapter{Video Image Processing}
\label{sec:video}

Video image acquisition is one of the most common sensory
devices used in robotics. Therefore \miro offers a common
{\tt VideoService} that can be adapted by the configuration file to any
supported video device.

Also, image processing is computationally expensive and many standard
filters exist in literature as well as some high performance
implementations such as \cite{IPP}. To facilitate the unified usage
within robotics scenarios, the video service provides a filter tree
framework.

Real-time image streams need a lot of bandwidth. This is usually
higher than the bandwidth available on most mainstream network
devices. Therefore the {\tt Video}-interface of \miro also provides
methods for usage of shared memory for sharing of images with its
clients on the same machine. Those methods are also designed for
minimizing copying overhead.

This section is organized as follows. First, the supported image
devices along with their bugs and features are discussed. Then the
filter tree framework is presented. First from an end user
perspective and second from a developer perspective. In section
\ref{SEC:VIDEO_INTERFACE} the video interface itself along with its
two use cases and helper classes is discussed.

\section{Video Device Access}

The VideoService currently supports the following devices for images
acquisition:

\subsection{Bttv Frame Grabbers}

These frame grabber cards are supported via video for linux \cite{}.
This is the standard way of connecting standard analog video cameras
to the computer.

Note, that the default number of frames used by the kernel driver for
video capture are two. In order to get the full frame rate (25/30Hz),
this has to be set to 4 buffers. This can be specified as a module
parameter.

\subsection{Firewire Digital Cameras}

\miro supports the fire wire digital camera protocol using
libraw1394\cite{} and libdc1394 \cite{}.
  
Note that most cameras 1394 controllers do not support being bus
master for dma transfers. Therefore the 1394 controller of the
computer has to be the highest numbered node. This can be configured
by a module parameter. Otherwise, unplugging the firewire cable from
the camera for some seconds helps.

\subsection{Matrox Meteor Frame Grabbers}

These rather old frame grabber cards are also supported by \miro.
Kernel drivers can be found at \cite{}. This device however, is mostly
unmaintained within \miro.

\section{Video Filter Trees}

For image processing many standard filters do exist that usually
perform a function $M_{n+1} = f_i(M_{n})$. Where $M_n$ denotes the input
image (matrix), $M_{n+1}$ the output image and $f_i$ is the actual
filtering function. 

Often filter chains are used (see figure \ref{fig:filterChain}). I.e.
first undistorting the image $f_0$ and then performing some color
indexing $f_2$. So the filtering function producing the output image
$M_o$ from the input image $M_i$ could be expressed as: $$M_o =
f_2(f_1(M_i))$$

\begin{figure}[!ht]
  \begin{center}
    \includegraphics[width=5cm]{filterchain}
  \end{center}
\caption{A simple filter chain.}
\label{fig:filterChain}
\end{figure}

More complex vision processing could also include the computation of
an edge image $M_e$ in addition to the color indexing above. This can
be done i.e. by producing a grey image $f_3$ and applying a canny
filter $f_4$, leading to the second filter chain: 
$$M_e = f_4(f_3(f_1(M_i)))$$

To avoid computing the undistorted image twice, a filter tree should
be used instead of two separate filter chains, as illustrated in
fiugre \ref{fig:filterTree}.

\begin{figure}[!ht]
  \begin{center}
    \includegraphics[width=6cm]{filtertree}
  \end{center}
\caption{A still quite simple filter tree.}
\label{fig:filterTree}
\end{figure}

As an additional complication, many sophisticated filters use more
then one input filter. For instance the Canny filter computes first
two images of the Zobel operator (one for the x axis $z_x$ and one for the y
axis $z_y$). So the above described filter tree is actually an acyclic
directed graph of filters, as illustrated in figure \ref{fig:filterGraph}.

\begin{figure}[!ht]
  \begin{center}
    \includegraphics[width=7cm]{filtergraph}
  \end{center}
\caption{A not that simple filter graph.}
\label{fig:filterGraph}
\end{figure}

\section{VideoService}

The program provided by \miro for image acquisition and preprocessing
is called {\tt VideoService}.

\miro comes with a set of basic filters. The term filter is used here
in a very general way, as almost every processing unit of the video
service is designed as a filter. Actually the image acquisition from
the different video devices is implemented as a filter, forming the
first set of filters available in \miro. The second set of filters are
basic image conversions, like byte order swapping, YUV to RGB
transformations etc. The third set of filters are the classic image
filters. Unluckily, there do not exist any of them currently in \miro.
It introduces further library dependencies for the middleware, which
we have to be careful about. And i.e. for Ipp based filters, we also
have a licensing problem. How to extend the video service to process
your own filters is covered in the following subsection.

Filter graphs and their individual parameters can be specified within
the configuration file of the robot. The VideoService actually does
not support every possible directed acyclic graph. The implications on
control flow and synchronisation are too cumbersome. Instead the
filters in the VideoService are organised as follows.

The filters are organised as a filter tree. That is, every filter has
one predecessor and $n$ successors. The filtertree is processed in
depth first order. This scheme is extended by so called filter
links: A filter can specify additional input filters. To guaranty the
correct order processing the filters, those additional filters have to
reside before the specifying filter, in terms of the depth first
evaluation scheme of the filter tree. Figure \ref{fig:legalIlligal}
tries to visualize this constrain.

\begin{figure}[!ht]
  \begin{center}
    \includegraphics[width=5cm]{filtergraphlegal}
    \includegraphics[width=5cm]{filtergraphillegal}
  \end{center}
\caption{A legal vs an illigal filter graph.}
\label{fig:legalIlligal}
\end{figure}

\subsection{VideoService Parameters}

The video service expects its configuration parameters with the
section Video of the configuration file. Its parameter is
also called Video. The
video service has the following parameters:

\begin{description}
\item[Width] The width of the input image (unsigned int).
  It has to be supported directly by the used video device.
\item[Height] The height of the input image (unsigned int).
  It has to be supported directly by the used video device.
\item[Palette] The palette of the input image (string).
  It has to be supported directly by the used video device.
  The recognized palettes of the video service are:
  \begin{description}
  \item[grey] 8 bit grey values.
  \item[grey16] 16 bit grey values.
  \item[rgb] 24 bit RGB format.
  \item[bgr] 24 bit BGR format (byte swapping).
  \item[rgba] 32 bit RGB format (plus alpha channel).
  \item[abgr] 32 bit BGR format (plus alpha channel).
  \item[yuv] 24 bit YUV format.
  \item[yuv411] compressed YUV format from 1394.
  \item[yuv422] compressed YUV format from 1394.
  \end{description}
\item[Filter] The root of the filter tree. The filter parameter
  has three entries:
  \begin{description}
  \item[Type] The filter type name (string).
  \item[Name] The name of the corresponding configuration file
    parameter entry (string).
    \item[Successor] The list of successor filters (Filter).
    \item[BackLink] The list of additional predecessor filters,
      specified by name (string). Note that this is the name of the
      filter instance, not the name of the filters video interface.
  \end{description}
\end{description}

\subsection{Video Filter Parameters}

For each filter in the filter tree, filter parameters can be given
within the parameter whose name matches the filter name given in the
filter tree. Those are also located within the section Video.

The base parameter set of each filter is as follows:

\begin{description}
\item[InterfaceInstance] Whether the filter has an instance of the
  video interface associated with it (bool).
  Each filter can have an instance of the video interface. That way
  every single node of the filter tree can be passed on to the client.
\item[Interface] The parameters of the interface. For each interface
  the following parameters can be specified:
  \begin{description}
  \item[Name] The name under which the interface is registered at the
    naming service (string). The default is Video. Note that if
    you have more than one Video object, you have to specify a unique
    name for each instance.
  \end{description}
\item[Buffers] The number of memory buffers, reserved by the buffer
  (unsigned int). The default is 4. While a client is accessing an
  image, the video service guarantees, that it does not become
  overridden with a new image from the video stream. If clients hold
  multiple images (i.e. for temporal integration) it has to be made
  sure, that there is still a buffer left, to put the next image from
  the video stream into.
\end{description}

\subsection{Video Device Parameters}

As mentioned above, all supported video devices are implemented as
filters. Therfore they inherit the base parameters. Note however, that
the device filters are not capable of having an interface
instance. Sharing memory mapped files is a bit tricky, and we didn't
want to go into that, yet. The common DeviceFilter parameters are:

\begin{description}
\item[Device] The fully qualified path of the video device
  (string). The individual device filters set this parameter to
  a common default (like /dev/video/0 for bttv and /dev/video1394/0 for
  firewire).
\end{description}

Common parameters for frame grabber based video devices are:

\begin{description}
\item[Format] The analog format of the video signal. This is either
  \texttt{pal, ntsc, secam} or \texttt{auto}.
\item[Source] The source of the video signal. Available sources are
  \texttt{composite1, composite2, composite3, svideo} and
  \texttt{tuner}.
\end{description}



The available video device filters are (listed by type name):

\begin{description}
\item[VideoDeviceBttv] The bttv device filter defines the following
  additional parameter:
  \begin{description}
  \item[Subfield] The subfield chosen (string). As the PAL and
    NTSC images are interleaved, the bttv frame grabbers can select,
    which of the half images to scan, if the height of the output
    image is half the height of the full image (384). Possible values
    are odd, even and all. Note however, that many frame grabber cards
    do not support this feature. In this case, use the FilterHalfImage
    described below.
  \end{description}
\item[VideoDevice1394] The firewire digital camera standard gives
  access to many of the internal camera parameters. These can be
  configured in the parameter section of the 1394 device. For all
  these values -1 denotes, that this parameter should be controlled by
  the camera.
  \begin{description}
  \item[Buffers] The number of image buffers used for dma transfers
    (unsigned int).
  \item[Brightness] (int).
  \item[Exposure] (int).
  \item[Focus] (int). Default is auto.
  \item[Framerate] (int). Default is 30Hz.
  \item[Gain] (int).
  \item[Gamma] (int).
  \item[Hue] (int).
  \item[Iris] (int).
  \item[Saturation] (int).
  \item[Sharpness] (int).
  \item[Shutter] (int).
  \item[Temperature] (int).
  \item[Trigger] (int).
  \item[WhiteBalance0] (int).
  \item[WhiteBalance1] (int).
  \end{description}
\item[VideDeviceMeteor] The matrox meteor device filter does not
  define any additional parameters.
\item[VideoDeviceDummy] A dummy device which can be used for offline
  testing.  It reads an image specified by the {\em Device} parameter.
  If the {\em Device} specifies a directory instead of a file, all
  {\em *.ppm} files of the directory are used by the device.
  \begin{description}
  \item[Timeout] The image refreshing time (sec).
  \item[Cyclic] (bool) The default is {\em true}. If true the device
    will restart with the first image after all images were loaded.
  \end{description}
\end{description}

\subsection{Basic Video Filters and Their Parameters}

\miro also provided the following basic filters with the standard
VideoService.

\begin{description}
\item[FilterCopy] As the video device filters cannot map the image
  data directly into shared memory for access by the clients, this
  filter is provided. It copies the image internally, to allow for an
  interface instance.
\item[FilterSwap3] Byte swapping for 24 bits per pixel images (BGR to RGB).
\item[FilterSwap4] Byte swapping for 32 bits per pixel images (ABRG to
  RGBA).
\item[FilterFlip] We have a camera that's mounted upside down. This
  filter flips the image. Also useful if your fly imitating robot has
  successfully landed on the ceiling.
\item[FilterHalfImage] This is a filter to extract a half image from an
  interlaced image as the bttv frame grabbers provide. It copies each
  second scan line. It defines the
  following additional parameter:
  \begin{description}
  \item[Odd] Take the second half image by starting with the second
    line (bool). Default is false.
  \end{description}
\end{description}

\subsection{Configuration Example}

Putting it all together the configuration section for the {\tt
  VideoService} might look like the following. This is actually the
configuration of our Performance PeopleBot. It has an analog video
camera and a bttv frame grabber card. The camera is mounted upside
down.

\begin{verbatim}
<!-- The video configuration section. -->
<section name="Video" >

 <!-- Parameter section of the VideoService -->
 <parameter name="Video" >

  <parameter value="bgr" name="Palette" />         <!-- Input image palette. -->
  <parameter value="384" name="Width" />           <!-- Input image width. -->
  <parameter value="288" name="Height" />          <!-- Input image height. -->
  <parameter name="Filter">                        <!-- Filter tree root. -->
   <parameter value="DeviceBTTV" name="Type" />    <!-- It's a bttv device. -->
   <parameter value="DeviceBTTV" name="Name" />    <!-- Params section name. -->   
   <parameter name="Successor" >                   <!-- Filter tree leafs. -->
    <parameter name="Filter" >
     <parameter value="FilterSwap3" name="Type" /> <!-- Byte swapping filter -->
     <parameter value="FilterSwap3" name="Name" />
     <parameter name="Filter" >                    <!-- Filter tree leafs. -->
      <parameter value="FilterFlip" name="Type" /> <!-- Upside down filter. -->
      <parameter value="FilterFlip" name="Name" />
     </parameter> 
    </parameter>
   </parameter>
  </parameter>
 </parameter>

 <!-- Parameter section of the bttv device. -->
 <parameter name="DeviceBTTV">                   
  <parameter value="/dev/video0" name="Device" /> <!-- Path of the device. -->
 </parameter>

 <!-- Parameter section of the byte swapping filter. -->
 <parameter name="FilterSwap3">
  <parameter value="true" name="InterfaceInstance" /> <!-- Video interface. -->
 </parameter>

 <!-- Parameter section of the upside down filter. -->
 <parameter name="FilterFlip">
  <parameter value="true" name="InterfaceInstance" /> <!-- Video interface. -->
  <parameter name="Interface">                        <!-- Interface params. -->
   <parameter value="Flipped" name="Name" />          <!-- Interface name. -->
  </parameter>
 </parameter> 

</section>
\end{verbatim}

\section{QtVideo}

{\tt QtVideo} is a test client for the {\tt VideoService} and the {\tt
  Video} interface. It displays an image stream from a {\tt Video}
interface instance and has an additional button, to save snapshots to
disk.

The {\tt QtVideo} tool accepts the following command line parameters:
\begin{description}
\item[-n] Name of the {\tt Video} interface instance within the CORBA
  naming service. Default is {\em Video}.
\item[-r] Remote access of the images. QtVideo will be using the
  methods for location transparent image access of the {\tt Video}
  interface. These are much slower, than the methods using shared
  memory buffers, but are the only option when running QtVideo on
  another machine.
\item[-v] Verbose mode.
\item[-?] Emitting command line help and exits.
\end{description}

\section{Video Interface}

The {\tt Video} interface is used to manage the access of image data
by client programs. It supports location transparent as well as
optimized local image access. Additionally connection management is
used, to switch off filter subtrees that are not accessed by any other
program.

The general access pattern of a client of a {\tt Video} object looks
like follows:
\begin{enumerate}
\item Get a {\tt Video} interface IOR from the naming service.
\item Connect to the {\tt Video} object.
\item Get images until done.
\item Disconnect from the {\tt Video} object.
\end{enumerate}

Note that due to the performance centric design of the Video
interface, the VideoService can easily be jammed by client programs
violating the access protocol of the interface. To facilitate the
correct usage, some helper classes are provided by \miro. These are
discussed in section \ref{SEC:C++HelperClasses}.

\subsection{Location Transparent Image Access}

To access an image as a client running on a different computer than
the {\tt VideoService}, the {\tt Video} interface offers the methods
\begin{itemize}
\item {\tt exportSubImage}
\item {\tt exportWaitSubImage}
\end{itemize}
They both return a copy of the image as return value. The first method
returns immediately the current image, while the second one waits until
a new image becomes available before returning. The image will be
scaled down to the size specified by method parameters.

Note, that due to the copying and network overhead of those methods,
they are only useful for debugging and monitoring purposes.

\subsection{Local Image Access}

For clients running on the same machine as the {\tt VideoService}, the
{\tt Video} interface offers the following methods for image access
via shared memory buffers:
\begin{itemize}
\item {\tt acquireCurrentImage}
\item {\tt acquireNextImage}
\item {\tt releaseImage}
\end{itemize}
While the first method returns immediately the buffer of the current
image, the second one waits until a new image becomes available before
returning. Note that the clients have to release each image buffer
after processing. Otherwise the VideoService will soon run out of
buffers, to share new images with its clients.

\subsection{C++ Helper Classes}

To facilitate the usage of the {\tt Video} interface and the adherence
of the connect/disconnect, acquire/release protocol, \miro provides
two simple helper classes. They are defined in the file {\tt
  \$(MIRO\_ROO)/src/miro/VideoHelper.h}

\begin{description}
\item[Miro::VideoConnection] This class connects to a {\tt Video}
  interface instance on construction and disconnects on destruction.
  The constructor takes a pointer to the {\tt Video} object, to
  connect to as argument.
\item[Miro::VideoAcquireImage] This class acquires an image from a
  {\tt Video} interface instance on construction and releases it on
  destruction. The constructor takes a reference to a {\tt
    Miro::VideoConnection} object, as first argument, the second
  selects whether the current or the next image is to be acquired.
\end{description}

\subsection{Example Video Client}

\lstinputlisting[frame=tb, first=12, caption={examples/video/VideoExample.cpp}]{VideoExample.cpp}
\label{lst:VideoExample}

\section{Video Broker Interface}

In addition to the Video interfaces of individual filters, the
VideoService also offers an interface that accesses the whole filter
tree of the service. It is called VideoBroker and also registers as
VideoBroker at the NamingService. It offers methods for synchronised
simultanious access to multiple filters and for inspection of the
filter tree.

\subsection{Synchronised Image Access}

The client side evaluation of the results of the various filters often
requires access to multiple filters, that originate from the same
input image. The following methods are available in the VideoBroker
interface for this purpose. They work quite like the corresponding
methods of the Video interface of the  individual filters.

\begin{itemize}
\item {\tt acquireNextImageSet}
\item {\tt releaseImageSet}
\end{itemize}

\subsection{Filter Tree Meta Information}

Statistical information about the filter tree can be optained at the
VideoBroker interface via the method:

\begin{itemize}
\item {\tt filterTreeStats}
\end{itemize}

It returns the number of connections, the processing time of the
filter as well as the processing time of the corresponding filter
subtree (not including the processing time of the filter), the name of
the filter, the ior and the name of its Video interface (if
available), the filter successors as well as the successor links.

\section{Writing Filters}

Sooner or later, it will become handy for a user to add 
custom filters to the video image processing framework. This is a two
step process. First the filter has to be implemented as a class
derived from the {\tt Filter} base class. Second a customized version
of the {\tt VideoService} has to be built. The later has two
reasons. The first is, that \miro currently does not contain a plug in
architecture for dynamically loadable modules. The second is, that
dynamic loadable modules require the usage of shared libraries. As
this is has some runtime performance impact, that might not really
desirable for performance critical tasks like image
processing. Nevertheless, patches are welcome.

\subsection{The Filter Base Class}

The base class of all filters is {\tt Video::Filter}. It provides
methods for filter tree setup and configuration, a harness for
recursive filter tree processing and buffer management in interaction
with a {\tt Video} interface instance.

The most important protected methods and data members for child classes are.
\begin{description}
\item[inputBuffer()] Returns a pointer to the input buffer.
\item[inputFormat\_] Specifies the data format of the input buffer.
\item[outputBuffer()] Returns a pointer to the output buffer.
\item[outputFormat\_] Specifies the data format of the output buffer.
\end{description}

\subsection{Methods to Overwrite}

The most important methods for child classes are.
\begin{description}
\item[Constructor] The constructor takes a struct of type {\tt
    Miro::ImageFormatIDL} as parameter. It describes the input format
  of the filter (hence the output format of the predecessor filter).
  The constructor of the {\tt Filter} base class copies this parameter
  into the two member variables {\tt inputFormat\_} and {\tt
    outputFormat\_}. The derived filter has to ensure, that the input
  format is a valid format for the filter and modify the member {\tt
    outputFormat\_}, to correctly describe the output format of the
  filter. If the input format is not valid a {\tt Miro::Exception} has
  to be thrown.
\item[process] The actual filtering method. Put your filter code here.
\end{description}

\subsection{Configuration and Parameter Processing}

If your filter needs additional parameters for configuration, your can
add them to the configuration file framework by specifying them in a
parameter file description, deriving them from {\tt
  Miro::FilterParameters}. You have to overwrite the factory methods
for {\tt Miro::FilterParamters} in the {\tt Filter} class to return an
instance of your derived parameters class. This can be automated by
using the macros:
\begin{description}
\item[FILTER\_PARAMETERS\_FACTORY(X)] For usage within the class definition.
\item[FILTER\_PARAMETERS\_FACTORY\_IMPL(X)] For usage within the class implementation.
\end{description}
Where X is the name of your derived Filter class. 

For initialization and cleanup the following two methods exist.
\begin{description}
\item[init] Called on initialization of the filter tree. It provides a
  pointer to the parameters object initialized from the configuration
  file. Use {\tt dynamic\_cast} to convert it to the type of your
  derived parameter class.
\item[fini] Called before destruction of the filter tree.
\end{description}

\subsection{Enabling and disabling features}

The interface also contains some flags that allow the
derived classes to specify the existens of certain properties for the
base class. Overwrite their settings in the constructor. Those are namely:
\begin{description}
\item[interfaceAllowed] Specifies, whether the filter result can be
  exported by a video interface. The default for Video::Filter is
  true. For Video::Device it is false (see section
  \ref{sec:VideoDevice} for details).
\item[inplace] Not implemented yet. (Allways false)
\end{description}

\subsection{Filter Meta Information: FilterImageParameters}

Sophisicated filters may also extract additional information from the
images, such as a region of interest etc. To hand this information to
subsequent filters, an additional interface does exist. It is based on
the same mechanisms as the filter parameters, except that they can't
be specified within configuration files. The base class is empty.

The factory methods also reside in the corresponding {\tt Filter}
class and are wrapped by the two macros:
\begin{description}
\item[IMAGE\_PARAMETERS\_FACTORY(X)] For usage within the class definition.
\item[IMAGE\_PARAMETERS\_FACTORY\_IMPL(X)] For usage within the class implementation.
\end{description}
Where X is the name of your derived Filter class. The base class is
named {\tt FilterImageParameters} and the macros expect the name of
any child class to end on {\tt ImageParameters}.

The {\tt Filter} interface provides accessor methods for the
parameters of the individual buffers:
\begin{description}
\item[inputBufferParameters] Returns a const pointer to the parameters
  of the input buffer.
\item[inputBufferLinkParameters] Returns a const pointer to the
  parameters of the linked input buffer specified by its index.
\item[outputBufferParameters] Returns a pointer to the parameters of
  the output buffer. Note that the parameter instances is not
  initialized with defaults before processing of the filter. That is,
  the filters {\tt process} method is responsible, that all data
  members of the image parameters instance contain meaningfull values
  on return.
\end{description}


\subsection{Example Video Filters}

To illustrate the framework, a simple example filter is provided under
{\tt \$(MIRO\_ROOT)/examples/videoFilter}.

\subsection{Calculating a Gray Image}

{\tt FilterGray} implements a simple gray
filter for images in 24 bit rgb format. The gray value is calculated
as the average of the weighted sum of the three color values. The
weighting can be specified by the filters parameters.

Start the example service in the directory videoFilter with \\
{\tt GrayVideoService -MiroConfigFile GrayVideoConfig.xml}

It registers two interfaces at the naming service: {\tt Video} and
{\tt Gray}. To examine the original and the filtered image, use {\tt
  QtVideo} and {\tt QtVideo -n Gray}.

\subsection{Image time series}

{\tt FilterDiff} illustrates the use of multiple input filters. It
calculates the difference of two input images. For that purpose it
locks images of one image source to provide a view into the past. The
number of images locked can also be specified by the filters
parameters. Note however, that it has to be less then the number of
output buffers provided by the input buffer.

Start the example service in the directory videoFilter with \\
{\tt GrayVideoService -MiroConfigFile DiffVideoConfig.xml} \\
Note that the usage of an input filter link is a bit artificial in
this configuration, as both input filters are the same.

It registers two interfaces at the naming service: {\tt Video} and
{\tt DiffImage}. To examine the original and the filtered image, use {\tt
  QtVideo} and {\tt QtVideo -n DiffImage}.

\section{Writing a new Input Device}

The root of a filter tree is a video device. A video device is
actually not an image filter but an image producer. But as mentioned
above, it is implemented as as specialization of the video filter
class. Some additional requirements and features do exist to handle
this case.

\subsection{BufferManger}

The video framework does use a buffer manager class to organize the
concurrent acces to the images of a filter by the filters and remote
clients. Video devices normally aint capable of having directly a
video interaface instance attached. The video interface uses shared
memory to allow for zero copy acces to the clients. Video devices
often use memory mapped io to acquire the images. Combining the two is
at least tricky, if possible at all. So most video devices just skip
it.

Instead the buffer manager class is used for synchronising the access
to the mmaped images between the filters and the hardware.



%%% Local Variables: 
%%% mode: latex
%%% TeX-master: "miro_manual"
%%% End: 

\chapter{Group Communication in Robot Teams}
\label{sec:NMC}

When it comes to multi-robot scenarios, group communication becomes an
issue of interest. As \miro is designed for usage in distributed
environments, most of the prerequisites for communication in teams of
robots are in place. But in the presence of WLANs, that have extremely
varying bandwidth, that dependents on the location of the computer
etc. some more requirements have to be met. Especially the TCP/IP
based transport protocol can become a bottleneck, when temporary
network breakdowns can regularly happen during operation.

\section{Event Channel Federation}

The event based communication of \miro is based upon the CORBA
Notification Service. In a team of robots, one central event
channel could be set up and used for team communication. While this is
most convenient from a user perspective, it creates a single point of
failure for a team of otherwise independent autonomous robots.

Using a single virtual event channel, that is distributed over
multiple channels, that run on the different robots helps both, the
programmers, as well as the systems fault tolerance. As such an event
channel federation is not part of the specification of the CORBA
Notification Service, \miro provides its own implementation. Namely
the Notify Multicast (NMC) module. Its implementation is based upon the work
done for the event channel federation of the RT-EC of TAO.

\section{Notify Multicast}

One notification service is run on each robot (computer). The NMC
module attaches one consumer and one supplier to the robots local
event channel. Additionally it subscribes at a multicast group. Events
that are subscribed by other robots are exchanged via the multicast
group.

For the automatic arbitration of the set of events that need to be
exchanged via the multicast group, the offer/subscription protocol of
the Notification Service is used. Offers to and subscriptions at the
event channel are determined by the \texttt{domain\_name} and
\texttt{type\_name} fields of the event header. In \miro the domain
name is by convention the name of the robot. The NMC module
posts all events that subscribed for at the local channel, but that
are not offered locally to the multicast group and collects these
posts from the other robots. Events that are requested by other
robots and offered locally, are posted to the multicast group on arrival
on the local event channel. Those events, that are subscribed for at
the local channel, but not offered locally, are published to the local
channel as they become posted on the multicast group.

\begin{figure}
%\vspace*{-0.5cm}
\begin{center}
\includegraphics[width=9.0cm]{federatedchannel}
\caption{\label{fig::federated_channel} 
A Federated Notification Channel Setup}
\end{center}
%\vspace*{-0.3cm}
\end{figure}

See also figure \ref{fig:federated_channel} for a possible
configuration of subscriptions and offers in a federated event channel.

\section{Usage}

The NMC module is part of most robot servers by default. It can be
enabled by the command line option -MiroNotifyMulticast or -MNMC for
short. It also exists, as a separate program: NotifyMulticast. 

\subsection{Parameters}

The parameters can be set in the robots configuration file. They are
located in the configuration group \texttt{Notification}. The
available parameters are.
\begin{description}
\item[MulticastGroup] The multicast group used for group communication.
\item[Timeout] The maximum age for events send over the group.
\item[Subscription] These events are shared with other robots event if
  no one seems to listen. - Events that are exchanged on the bases of the
  offer/subscription management protocol need some time, until the
  request is acknowledged by the suppliers. By providing events with
  this option, this time can be skipped. Note, that only the type name
  is specified here. The domain name is assumed to be the robots name.
\end{description}


%%% Local Variables: 
%%% mode: latex
%%% TeX-master: "miro_manual"
%%% End: 

\chapter{Logging}
\label{sec:logging}

The logging of debug, information and error messages can help a great
deal in the recovery from failure situations. This is especially true
for robotic applications. Unfortunately the sheer mass of console
output from the different modules can become too complex quite fast. So
\miro comes with a set of logging facilities that cover different
levels of the system functionality. The first framework, that is
introduced within this chapter is based on the logging framework
provided by ACE. It is a system log oriented framework to organize printf
like output on various levels. Chapter \ref{sec:ecLogging} will cover
another, more sophisticated set of logging functionality.

\section{Log Levels and Categories}

\miro uses two mechanisms to customize logging output.  Log levels are
used to give fine grained control over the verbosity of log output.
\miro defines log levels from 0 to 9. The higher log levels are
considered to be used for debugging purposes only. They are therefore
also referred to as debug levels. They usually are only meaningful for
developers. For debug levels log categories 
are used to restrict the logging output to one or more subsystems of
the robot application. Output from debug levels is only displayed if
its log category is enabled, too.
Additional log categories can be defined for your own
programs.

The different log levels are:
\begin{enumerate}
\setcounter{enumi}{-1}
\item \emph{Emergency} Log level of messages reporting an emergency.
  Your robot is on fire etc. This log level is not maskable, except if
  you turn of logging at configure time.
\item \emph{Alert} Log level of messages reporting an alert. This is
  when the red lights start blinking and this unnerving honking sound
  is played.
\item \emph{Critical} Log level of messages reporting a critical
  condition. This usually is an unrecoverable error, that leads to
  the termination of reporting program.
\item \emph{Error} Log level of messages reporting an error.
  This indicates a real error, but the program will usually try to
  continue anyway.
\item \emph{Warning} Log level of messages reporting a warning.  A
  warning should be fixed, but the program is likely to work anyways.
\item \emph{Notice} Log level of messages reporting a notice. Make
  a post-it and add it to the other 500 ones.
\item \emph{Ctor\_Dtor} Debug level of messages reporting a
  constructor/destructor entry.  This debug level is designed to hunt
  segmentation faults on startup and exit. --- This is when all the big
  ctors/dtors are run.
\end{enumerate}

The debug levels are:
\begin{enumerate}
\setcounter{enumi}{6}
\item \emph{Debug} Log level of messages reporting debug output.
\item \emph{Trace} Log level of messages reporting program trace
  output.
\item \emph{Prattle} Log level of messages reporting really verbose
  comments on the program execution.
\end{enumerate}

The different log categories used within \miro are the
following. (All brand names have to be marked as such, as soon as we
find the time):
\begin{description}
\item[Miro] This category is used by the \miro core components,
  namely those located in the library libmiro.so
\item[Video] Used by the components of the video image processing
  system. See section \ref{sec:VIP} for details.
\item[NMC] The notify multicast based event channel deliberation, as
  used for team communication. See section \ref{sec:NMC} for details.
\item[B21] Used by the components exclusive to the B12 robots of RWI.
\item[Pioneer] Used by the components exclusive to the Pioneer robots
  or Active Media.
\item[Sparrows] The custom built Sparrows platform.
\item[Faul] The motor controllers from Faulhaber.
\item[DP] The pan-tilt unit from DirectPerception.
\item[Sick] The sick laser scanners have their own category.
\item[Sphinx] The sphinx speech components.
\item[Dtlk] The DoubleTalk speech components.
\item[BAP] The reactive control subsystem logs messages within this
  category. See section \ref{sec:BAP} for details.
\end{description}

\section{Run-Time Configurability}

The log level and log categories that are displayed can be configured
at program startup by command line parameters.

\begin{description}
\item[-MiroLogLevel $<$n$>$] or -MLL $<$n$>$ for short, selects the log level
  up to which data is logged. The log levels are referenced by number.
  The default log level is 4. Log level 0 can not be masked by command
  line parameters, but only by turning off logging entirely at compile
  time.
\item[-MiroLogFilter $<$name$>$] or -MLF $<$name$>$ for short, selects a
  category to log. This option can be specified multiple times to
  enable multiple categories. The category enabled by default is
  \miro.
\end{description}

\section{Usage in Source Code}

This logging facility is used extensively in the \miro sources. It can
also be used and extended for the usage in robotic projects based on
\miro.

\subsection{Miro::Log}

The class \texttt{Miro::Log} is defined in \texttt{miro/Log.h}. It
works mostly as a namespace for the logging facility. It holds
constants for all log levels as well as log categories. The main
method is \texttt{Miro::Log::init()}. This is a static method that
takes \textit{argc} and \textit{argv} as arguments, to parse for the
\texttt{-MiroLogLevel} and \texttt{-MiroLogFilter} command line
options. Additionally, accessor methods and predicates are defined to
query log levels etc.  at runtime.

\subsection{Macros}

\texttt{miro/Log.h} also defines a set of macros that are usable for
development with, and within \miro.
\begin{description}
\item[MIRO\_LOG] Produce some log output, if the required log
  level is enabled. It takes two arguments. The first is the log level,
  the second is a single character string. The log level is specified
  by name. The canonical name format for log levels is the log level
  name in capitals, prefixed by \texttt{LL\_}.
\item[MIRO\_LOG\_OSTR] Like \texttt{MIRO\_LOG}, but the second
  parameter is used as right hand side argument for an output stream
  operator $\ll$. So it can contain some operator $\ll$ concatenated
  expression.
\item[MIRO\_DBG] Produce some log output, if the required log level as
  well as the log category is enabled. It takes three arguments. The
  first is the log category, the second the level, and the third is a
  single character string. The canonical name format for log
  categories is the log category name in capitals.
\item[MIRO\_DEBUG\_OSTR] Like \texttt{MIRO\_DBG}, but the last
  parameter is used as right hand side argument for an output stream
  operator $\ll$. So it can contain some operator $\ll$ concatenated
  expression.
\item[MIRO\_LOG\_CTOR] For constructor tracing. Accepts a single
  parameter containing the class name as string. Logged with log level
  6.
\item[MIRO\_LOG\_DTOR] For destructor tracing. Accepts a single
  parameter containing the class name as string. Logged with log level
  6.
\item[MIRO\_DBG\_TRACE] For method call traces. Takes the log
  category as argument.
\item[MIRO\_ASSERT] The standard assert macro. It is provided to
  enable disabling of assert macros in (inline) code of \miro without
  disabling them for user code too.
\end{description}

\section{Compile-Time Configurability}

Log and debug messages can be disabled entirely at compile time by two
configure flags: 
\begin{description}
\item[--disable-DebugInfo] disables debug information, that is namely
  log levels above 6. The \texttt{MIRO\_DBG\_...} macros will be
  replaced by no-op implementations, removing footprint and performance
  overhead of the programs, introduced by debug information. -
  Especially the performance overhead is mostly negligible, so this
  option might well be omitted even for release versions.
\item[--disable-LogInfo] disables log information entirely. Like
  above, this can save some footprint but is hardly measurable in
  performance.
\end{description}

\section{Test and Example Programs}

The programs located at tests/log provide some testing facilities
for the logging facility and can serve as a practical example on the
setup of the logging facility as well as on the usage of
the logging and debug macros. The only program currently located there
is TestLog.cpp. The source code is documented to help understanding.

%%% Local Variables: 
%%% mode: latex
%%% TeX-master: "miro_manual"
%%% End: 

\chapter{Event Channel logging}
\label{sec:ecLogging}

The acquisition of data during the run of a robot application from
various levels of sensor processing is an essential feature for
debugging, evaluation and performance assessment. The event based
communication paradigm of \miro is designed to distribute raw sensor
data as well as higher level system events, like the latest belief
state of the robot. The stream of events resembles therefore a quite
complete trace of the system state during the run of a robot program.
\miro provides functionalitiy to log such an event stream generically
to a file. The data can not only be reread into the system, but the
events can also be redistributed over the event channel, allowing for
detailed offline analysis. In \cite{Utz+Mayer+Kraetzschmar:04} the
various configuration and application scenarios are described along
with a detailed performance analysis of the facility. This chapter is
concerned with a more technical view, like the setup of the logging
client, the replay as well as the maintenance of logged data and the
file format.

\section{Logging Client Class}

\section{Parameters}

\section{Standalone Clients}

\subsection{Command Line Parameters}

\subsection{Event Channel Client}

\subsection{NMC Client}

\section{LogPlayer}

\section{File Format}

\section{Test Programs}

%%% Local Variables: 
%%% mode: latex
%%% TeX-master: "miro_manual"
%%% End: 

\chapter{Behaviour Engine}
\label{SEC:BEHAVIOURS}

Controlling the actuators of an autonomous mobile robot is one of the
central aspects of mobile robot research. As \miro exposes the
interfaces to the motor controllers etc. of the mobile platform it
enables researchers to easily evaluate new approaches to model sensor
actor control loops. Be aware that at this point the latencies
introduced by various levels of the robot architecture can become a
critical issue. To eliminate the network latency and prevent yourself
from occasional network bandwidth problems we recommend that you run
your control programs collocated on the same computer with the service
that accesses the actuators device.

\miro supports the behavioural control paradigm introduced by Brooks
\cite{Brooks} by its own behaviour engine. It is designed to allow for
a quick start into behaviour robotics writing your own behaviours. Yet,
due to its open and extensible design it is also capable of handling
sophisticated control tasks as demonstrated by its use within the
RoboCup-scenario by our \sparrow robots.

\section{The Concept of Behaviours}

The basic idea of the behaviour approach to robot control is as
follows. Instead of the sense-plan-act paradigm of classical AI, the
task is splited into a set of reactive behaviours, that each try to
fulfill a small subtask of the problem set. For each of the tasks only
a very limited part of world modeling is needed (often, even raw
sensor readings are sufficient). By combining the output of the
various behaviours by an arbiter, the emergent higher level behaviour
of the system is achieved, solving the recommended high level task.

An accepted bottleneck of this approach is the arbitration and
calibration of large sets of behaviours necessary to fulfill different
aspects of a high level task. Therefore in the \miro framework the
behaviour engine also support the hierarchical decomposition of
behaviour sets by allowing them to be grouped in so called action
patterns that can be activated alternately by so called transition
messages. A set of action patterns is called a policy within the \miro
framework. 

There do exist two different kinds of transitions within the behaviour
framework: local and global transitions. Local transitions are emitted
by the behaviours by name. For example a transition named
``GoalReached'' might be emitted by a behaviour when it decides, it
has fulfilled its task. Within an action pattern a transition is
linked to the successor pattern, that will be activated as the
transition message occurs. That way totally different action patterns
can be activated, when the behaviour emits the ``GoalReached''
transition message within different action patterns. Global
transitions are somehow simpler. They contain directly the action
pattern, that is to be activated next. This mechanism is designed for
modules of the robot, that are external to the behaviour framework.
For example if a task planner decided it would be time for the robot
to go home, it might activate the ``GoHome'' pattern just by issuing a
global ``GoHome'' transition. The robot will then switch to its GoHome
pattern regardless of what it was doing before.

\section{Introductory Examples}

Lets look at the simple action pattern, who's high level task is to
explore the environment by performing a random walk. This is easily
splited into two distinct behaviours. The first subtask is not to
collide with the environment. This can be achieved by an avoid
behaviour, that reads the front sonar sensors to determine how far it
is away from the nearest obstacle. If the minimal distance is below
some threshold, it tells the arbiter to turn away into another
direction. The second behaviour would be a wander behaviour. It selects
from time to time just randomly some translational and rotational
velocity, making the robot move around. It does not have to care about
obstacles, since those are taken care of by the avoid behaviour.

The task of choosing the actual velocities to be applied to the motors
is performed by the arbiter. It therefore plays a central role in the
behaviour approach. There exist various kinds of arbiters, all
choosing different policies for this task. But a simple priority based
arbiter suffices for many scenarios. In \miro there currently just
exists a priority based arbiter, but since it is an extensible
framework you can easily plug in your own one. The priorities in this
example could be applied straight forward. The avoid behaviour has
higher priority as the wander behaviour.

Note the easy extensibility of this approach. For example, if we have
some bump sensors that indicate collisions with the environment. We
could just add another, let's call it ``EmergencyStop'' behaviour and
assign it the highest priority.  If one of the bumpers is pressed it
makes the robot stop and wait for rescue by one of the operators.

As a further extension we assume we have two tasks. The first is the
random walk described above, the second would be a wall following
behaviour, that simply is capable to drive the robot in a defined
distance along a wall. The decision what to do is provided by an
external source, say, a button located on the robot and pressed by the
operator when demonstrating the management the capabilities of the
newly bought autonomous mobile robot. The second action pattern looks
quite like the first one, except that the wander behaviour is
exchanged by a wall following behaviour, that drives the robot along
the wall. (Note that this behaviour does not have to care about walls
thar are blocking the way at the end of the corridor, since this can
be taken care of by the avoid behaviour.) These two action patterns
now form a policy. The pressing of the button sends a transition
message, that disables the currently running action pattern and
enables the other one.

The action pattern / transition message mechanism also fits naturally
for coupling reactive behavioural control with deliberative planning
architectures. The transition could also be raised by a path planner
and the corresponding action patterns could be 'move to point', 'drive
through door' or 'dock at power supply'. Indeed this was already done
within \cite{Hans-Diplom}.


\section{Example Usage}

As an example of the behaviour framework in use,
let us start the obstacle avoid demo.
\begin{itemize}
\item
  Start all the robot main services from within {\tt \$MIRO\_ROOT/bin/}
  with: \\
  {\tt PioneerBase} \\
  This server has to run on the robot.
\item
  Start the Behaviours, specifying the policy at the command line: \\
  {\tt Behaviours ../examples/behaviours/PpbStraight.xml} \\
  This client can run on any computer, however,
  it is a good idea to run the behaviours on the robot,
  to avoid that network latencies stop the control loop.
\item
  Start the policy controller GUI from within
  {\tt \$MIRO\_ROOT/bin/} with: \\
  {\tt PolicyController} \\
  to start and stop the policy run by Behaviours.
  Run this from any computer.
\end{itemize}

\subsection{The Behaviour Control Loops}

The core of the behaviour framework gets instantiated and initialized
in the file {\tt Behaviours.cpp} in {\tt
  \$MIRO\_ROOT/examples/behaviours/} with its {\tt main()} procedure.
It instantiates all available behaviours and arbiters starts a thread
for the timer service (ACE\_Reactor) needed by timed behaviours (see
section \ref{SEC:MIRO_TIMED_BEHAVIOUR}) and runs the CORBA event loop
for its BehaviourEngine interface, as well as for the event consumers
like event behaviours (see section \ref{SEC:MIRO_EVENT_BEHAVIOUR})
that get pushed by {\tt PioneerBase} for instance.  Behaviours
themselves do not have an own {\tt main} procedure.  They are
implemented in the {\tt action()} method of ordinary C++ objects which
have to be known in {\tt Behaviours.cpp}.  Necessarily, all behaviours
run on one computer.

\begin{figure}[!ht]
  \ifpdf
  \includegraphics[width=\textwidth]{../fig/Miro_behaviours.png}
  \else
  \includegraphics[width=\textwidth]{../fig/Miro_behaviours.eps}
  \fi
  \caption{The Behaviour Architecture Overview}
  \label{FIG_BEHAVIOUR_ARCHITECTURE}
\end{figure}


As seen in Fig.~\ref{FIG_BEHAVIOUR_ARCHITECTURE}, {\tt Behaviours.cpp}
runs two loops synchronously: the timer loop for timed behaviours and
the CORBA loop.
The timer loop is started in a separate thread (invoked by the command
"task$\rightarrow$open(NULL);").
It schedules the timed behaviours
by invoking their {\tt action()} methods one after another. 

The CORBA loop of {\tt Behaviours.cpp} (invoked by the command
"server.run(5);").  processes requests for the BehaviourEngine
interface (which it registers under the name "BehaviourEngine" at the
CORBA Naming Service) from the PolicyController, as well as it
administers all CORBA requests from {\tt PioneerBase} and invokes the
event behaviours.  The {\tt action()} method of an event behaviour is
invoked whenever the CORBA notification service publishes an event to
which the behaviour has subscribed.  Events are subscribed to by their
names, e.g. "Tactile" or "Sonar" (see {\tt TactileStop.cpp} or 
{\tt SonarNotify.cpp}). The published sensory events and their associated
payload are described in section \ref{SEC:ASYNCHRONOUS_SENSORY_INFORMATION}.

Invocation of an {\tt action()} method by the event "Tactile"
is reminiscent of a call to a server registered as "Tactile" by a client,
e.g.\ {\tt PioneerBase}. However, there are a few differences:
{\sl (i)} the event "Tactile" calls not a program with its own
{\tt main()}, but only an {\tt action()} method,
{\sl (ii)} not one unique method is invoked,
but all {\tt action()} methods which have subscribed to that event and
{\sl (iii)} unlike proper services, the TAO {\tt nslist} utility does
not list these event consumers,
because they cannot be meaningfully invoked by an external program.


The Behaviour Engine Implementation loads an XML Policy File into a
policy instance and starts and stops the policy. On activation of an
action pattern, all behaviour that are part of the action pattern get
connected to the event chanel or the timer respectively.

\section{Arbiters and Messages}

Arbiters select the action from all the suggestions the behaviours of
an action pattern make. The behaviours make a suggestions by sending
an arbitration message to the arbiter. Each arbiter has its own
arbitration message format, depending on the arbitration scheme used
and the actuatory device(s) it controls.

In \miro currently exist three arbiters. The {\tt MotionArbiter} used
in all examples, the specialized {\tt PowerArbiter}, that uses the
SparrowMotion interface and the experimental {\tt WindowArbiter}.

\subsection{MotionArbiter}

The arbiter acts only on the motor output of the behaviours:
instead of sending motor commands directly to the {\tt Motion} interface,
the behaviours invoke the method arbitrate() of the arbiter,
passing it a message.

The {\tt MotionArbiter} uses a simple fixed priority, winner
takes all, arbitration scheme. That is, each behaviour has a fixed
priority that corresponds with the order the behaviours are defined
within the action pattern. In the {\tt MotionArbiterMessage}, that is,
the parameter passed by the {\tt arbitrate()} call, the behaviour can
specify whether it wants to suggest the arbiter an action or not. The
action is defined by the translational and rotational velocity for the
motion. The arbiter then decides whether to transform the content of
the message into motor commands or whether to issue the motor commands
of another, conflicting behaviour which has a higher priority.

Note that the arbitration of messages before conversion into motor
commands, however, is only for the wheels of the robot by the {\tt
  MotionArbiter}.  The other actuators, for instance the camera
pan-tilt unit or the grippers, are to be controlled directly in this
case. No arbiter has to be defined for them, as long as a only single
behaviour is expected to control them.

\section{Implementing a Behaviour}

All behaviours are derived from the base class {\tt Miro::Behaviour}.
The central method of this class is the method {\tt action()}.
It is to be overwritten by the programmer
and has to contain the behavioural code.
The method becomes invoked by the behaviour control file,
i.e.\ {\tt Behaviours.cpp}, which includes the behaviour header files
and which links the behaviours to their corresponding control loops.
However, if only a new policy is combined from existing behaviours,
then {\tt Behaviours.cpp} does not need to be modified.

The behaviours' {\tt action()} method is expected to call an arbitration
method or send a transition message. Note the inversion of control flow.
A behaviour is not allowed to jump into some infinite loop,
but the {\tt action()} method is called continously
as long as the behaviour is active within some action pattern.
So every {\tt action()} method must return quickly enough after its
invocation to give time to the others. This implements the idea of
simple reactive behaviours rather than complex strategies.
If it is necessary to run a behaviour continuously in parallel to other
behaviours, then own threads have to be created, e.g. using ACE.

Only one instance of each behaviour object is used within the behaviour
engine, even for different action patterns.  Therefore data that needs
to be stored between subsequent {\tt action()} method invocations can
be easily stored within member variables.

From an implementation perspective, there are three kinds of
behaviours, depending on how the data, the behaviour bases its
decisions on, is delivered.

\subsection{Miro::TimedBehaviour}
\label{SEC:MIRO_TIMED_BEHAVIOUR}

This is the base class for a timer scheduled behaviour. A behaviour
derived from this class runs with all its brothers and sisters
cooperatively multi-threaded in one thread of control. The pace at
which its {\tt action()} method is called is selectable by a parameter of
the base class.

This is the most simple form of behaviour design and in many cases most
straight forward. It is especially suitable for behaviours that poll
their sensory information of the world model, like for instance the
current scan of the laser range finder for collision avoidance. Also
the above mentioned wander behaviour - which doesn't use any sensory
information would be best implemented as a child of the
{\tt Miro::TimedBehaviour} class.

See {\tt \$MIRO\_ROOT/examples/behaviours/simple/Wander.cpp|h} for a
complete source code example.

\subsection{Miro::EventBehaviour}
\label{SEC:MIRO_EVENT_BEHAVIOUR}

This is the base class behaviours using asynchronous sensory
information published by the Notification Service. It subscribes for
the events it likes to get pushed and the action method is called
whenever a new message arrives. A pointer to the current structured
event is then available as member variable.

The emergency stop behaviour described
above would be a good candidate for such a behaviour, since it would
have to poll excessively as a timed behaviour in order to minimize the
latency between a bumper pressing a the actual stop, while it would
only need to call an arbitration method in very rare occasions.

See {\tt \$MIRO\_ROOT/examples/behaviours/simple/TactileStop.cpp|h} for a
complete source code example.

\subsection{Miro::TaskBehaviour}
\label{SEC:MIRO_TASK_BEHAVIOUR}

Behaviours derived from this calls run within their own thread of
control, not blocking others even if they need fairly long for their
decisions. Note that such a behaviour is likely to be miss designed,
since it contradicts the behavioural approach to need excessive time
to come to a decision. But if you have need for something like this,
since you are doing something we didn't think of, this is the class to
base your behaviour on. Note however, that since behaviours are shut
down cooperatively, also a task behaviour is not allowed to loop
indefinitely within its action method. If the behaviour is still part
of the currently running action pattern, its action method will be
call immediately after giving control back to the behaviour framework.

\subsection{Behaviour Parameters}

Behaviours usually get used within different action patterns. But they
are often expected to behave slightly different within each
constellation. Therefore each behaviour has an associated {\tt
  BehaviourParameters} class which is designed to hold the different
parameter sets for the different use cases of a behaviour. These
parameter classes become initialized on startup of the behaviour
framework and are expected to be static and constant during the run of
an entire policy. (How to handle dynamic parameters like destination
coordinates of a 'move to position' behaviour is explained in section
\ref{SEC:DYNAMIC_PARAMETERS})

Behaviours, which subclass the {\tt BehaviourParameters} class,
need to have their own {\tt BehaviourParameters}
subclass factory methods, for dynamic instantiation. This is captured by
the two macros:
\begin{itemize}
\item BEHAVIOUR\_PARAMETERS\_FACTORY(X)
\item BEHAVIOUR\_PARAMETERS\_FACTORY\_IMPL(Y, X)
\end{itemize}
X denotes the name of the behaviour parameters class and Y denotes the
name of the behaviour class. The first macro has to be placed within
the class definition and the second within the cpp-file containing the
behaviour implementation.

The behaviour parameters classes are handled by the parameters
framework described in the previous chapter. Section
\ref{sec:BehaviourDescription} covers the details for parameter
classes for behaviours.

\subsection{Behaviour Initialization}
\label{sec:init}

Before an action pattern becomes activated all its behaviours {\tt
  init()} methods are called successively. This allows behaviours to
initialize their per task parameters (like destination coordinates) in
a convenient way. Note however, that the {\tt init()} method can be
called while the behaviour is already active (see section
\ref{sec:active-inactive}) and so its {\tt action()} method can be
concurrently running. Therefore a mutex is needed to avoid race
conditions. For an example see:

{\tt \$MIRO\_ROOT/examples/behaviours/simple/Timer.cpp|h}

Parameters that are valid for the whole lifetime of a behaviour, such
as references to the robots services or other objects within the
behaviours address space are best to be passed during construction
of the behaviour, forming a so called initializing constructor.

\subsection{Behaviour Activation and Deactivation}
\label{sec:active-inactive}

When a behaviour is to become active due to it being part of an action
pattern its {\tt open()} method is called. If the behaviour is no longer
part of the next to be running action pattern its {\tt close()} method is
called. If a transition from one action pattern to another is
performed and the behaviour is part of both behaviour sets, then no
calls to {\tt close()} and {\tt open()} methods will be issued. Only the
behaviours {\tt init()} method will be called, to allow it to update its
parameter sets. This is useful to avoid unnecessary behaviour
shutdown.

Note also, that the {\tt close()} method can be called while the behaviour
is concurrently within the {\tt action()} method. --- Its call to the
arbiter will then just be ignored. On the call of {\tt open()} however it
is guaranteed that the behaviour currently is not running.

\subsection{Changing Behaviour Parameters within an Action Pattern}
\label{SEC:DYNAMIC_PARAMETERS}

Another way for interaction with the behaviour framework from outside
is by changing the parameters of a behaviour within an action pattern
during execution of the behaviour engine. In this case the behaviour
name as well as the action pattern it belongs to has to be known. You
can then query as well as set the behaviour parameters instance via
the interface of the policy.

However there are some issues that have to be take care of. First, the
type of the derived behaviour parameters struct has to be known, as
you get a pointer to a BehaviourParameters instance, that has to be
down casted appropriately. Second, to avoid memory leaks and
segmentation faults, the memory handling has to be understood. When
querying the parameters, the caller takes ownership of the object
returned. When setting the parameters, the action pattern takes
ownership of the passed parameters instance. - Therefore the object
has to be allocated on the heap (that is, with {\tt new}). But
actually that's not too much of a problem, if you stick to the
intended protocol: 
If you first retrieve the parameters instance by querying the policy,
change the parameters and write them back by the set method,
everything works okay.

A similar interface exists for the behaviours themselves. But they only
can query and set the parameters of behaviours in action patterns that
are linked to the current action pattern by a local transition. The
behaviours can therefore only query the parameters by the transition
name instead of the action pattern name.

\section{Arbiters}

The arbiter framework is similar to the behaviour classes. It will be
explained in more detail as soon as we do some more work on arbiters.

\section{Building Action Patterns}

``Now I built all my behaviours. What code do I have to write to make
them an action pattern?'' Well, you don't have to write code. Action
patterns and policies are defined within an XML file. Allowing for
fast and convenient modifications. Which is especially cool when
debugging. An Instance of each arbiter and behaviour has to be
registered within repositories. Afterwards, the policy can configure
the action patterns and transitions by parsing the XML file. 
To make live easier, the editing of action patterns, transitions and
behaviour parameters can also be done using the GUI based PolicyEditor
(see section \ref{SEC:POLICY_EDITOR}), but
let us first take a look at how those things are put together in the
XML syntax.

\subsection{The Policy File}

A policy file can have the following tags and attributes:

$<$!--MiroPolicyDocument--$>$

\begin{description}
  \item[policy]
    One file always describes one entire policy. A policy defines a
    state machine that consists of a set of action patterns. Each
    action pattern describing one state of the machine.
  
  \begin{description}
  \item[actionpattern]
    An action pattern describes a set of behaviours, that can run
    simultaneously and produce some emergent higher level behaviour.
    

    Attributes:
    \begin{itemize}
    \item \textit{name} The name of the action pattern.  
    \item \textit{start (true/false)} The action pattern to be active
      at startup is to be marked true, the others false or unmarked.
    \item \textit{x, y} For GUI purpose, just ignore them.
    \end{itemize}
  
    \begin{description}
    \item[behaviour]
      That's what all the chapter is about. Read it again.

      Attributes: 
      \begin{itemize}
      \item \textit{name} Name of the behaviour.
      \end{itemize}  

      \begin{description}
      \item[parameter]
        Behaviours parameters can be specified within the policy file
        for each action pattern. 

        Attributes: 
        \begin{itemize}
        \item \textit{name} Name of the parameter.  
        \item \textit{value} The value of the parameter.
        \end{itemize}
      
      \end{description}  

    \item[transition]
      Message that triggers a transition to another
      action pattern.  

      Attributes: 
      \begin{itemize}
      \item \textit{target} The action pattern to activate next.
      \end{itemize}  

    \item[arbiter]
      To decide which action to choose from the different outputs of
      the different behaviours one needs to arbitrate one way or the other.

      Attributes: 
      \begin{itemize}
      \item \textit{name} The name of the arbiter to use.
      \end{itemize}
    \end{description}
  \end{description}
\end{description}

\subsection{The Repositories}

A policy can be built on the basis of an XML description. For the
parser of this description to be able to construct the policy, it has
to be able to refer to instances of behaviours and arbiters that are
mentioned within the XML file by their name.  For this purpose a {\tt
  Miro::BehaviourRepository} and a {\tt Miro::ArbiterRepository} class
do exist.  At these Repositories an instance of each behaviour and
arbiter has to be registered. The name of the behaviour has to be
reported by the {\tt behaviourName()} method.  Since we only need one
instance of each of these repositories, there do exist a global
instance of each. A pointer to such an instance can be obtained by
the classes static method {\tt instance()}.

\subsection{The Behaviour Factory}

All the initialization stuff necessary before constructing an instance
of the Policy class can be done within the main function of your
program. This task consists mostly of obtaining the needed object
references, instancing
behaviours and arbiters and registering them at their respective
repositories.  This can be bundled within a so called behaviour factory
class as can be seen in the behaviour example at:

{\tt \$MIRO\_ROOT/examples/behaviours/simple/BehaviourFactory.cpp|h}

It uses the simple base class:

{\tt \$MIRO\_ROOT/src/miro/BAFactory.cpp|h}

\subsection{The Behaviour Engine}

There does exist a small CORBA based interface for interaction with
the behaviour framework. It is the {\tt BehaviourEngine} interface and
the {\tt PolicyController} (section \ref{SEC:POLICY_CONTROLLER}) is a
simple GUI based client for this interface. To allow your own
behaviour sets to be controlled via this interface, the policy
controller implementation has to be instantiated within your
executable. This can be seen in the example at:

{\tt \$MIRO\_ROOT/examples/behaviours/engine/SimpleBehaviourEngine.cpp}

\section{The Policy Editor}
\label{SEC:POLICY_EDITOR}

Editing large XML files is tedious and error prone. Therefore \miro
offers an GUI based editor, with which you can build and edit your
policies. You can also edit the parameters defined within the
Parameters classes of your behaviours.

The policy editor can be started with no parameters or with the policy
file to edit as first parameter. It has a menu bar and displays the
policy graph within the big scroll view area of the program.

\subsection{The Menu}

\subsubsection{File}
\begin{itemize}
\item New \\
  Create a new policy.
\item Open ... \\
  Open a policy file. The file dialog will be displayed.
\item Save \\
  Save the current policy to disc. If no file name is yet specified,
  Save as will be invoked instead.
\item Save as ... \\
  Save the current policy under a new file name. The file dialog will be
  displayed.
\item Send to ... \\
  Send the current policy to a robot. A dialog box will be displayed to
  enter the robots name (the naming context, its behaviour engine is
  registered at). The behaviours and arbiters within the policy have to
  be present within the behaviour and arbiter repositories in the
  behaviour engine running on the robot. - Otherwise the parsing will
  fail.
\item Quit \\
  Quit the policy editor.
\end{itemize}
\subsubsection{Options}
\begin{itemize}
\item Behaviour Descriptions \\
  Display the behaviour descriptions dialog. The dialog consist of a
  listbox with the currently loaded behaviour parameter description
  files. (See the chapter \ref{SEC:PARAMETERS} for more details on
  parameter files.) Buttons for adding and deleting files are
  provided. On pressing the button add, the file dialog will be
  displayed. Note that the selected file has to contain a valid
  behaviour description. - Verifying, whether the load was successful
  is currently only done on leaving the dialog. On the positive side
  the currently loaded description is memorized in the file
  .PolicyEditorConfig.xml file in the users home directory for further
  runs of the policy editor.
\end{itemize}
\subsubsection{Help}
Currently this menu only contains two about dialogs.

\subsection{Editing the Policy Graph}

Clicking with the right mouse button into the editing area produces
a popup menu with the option to place a new action pattern. Enter the
name of the new pattern in the dialog box.

Clicking with the right mouse button into the name of an action
pattern produces a popup menu with the following options.
\begin{itemize}
\item Start pattern \\
  Select the pattern to be the start pattern. That is the pattern, the
  policy will begin with by default. Every policy has to have exactly
  one start pattern. The start pattern is marked with two asterix
  around the pattern name.
\item Add behaviour \\
  Add a behaviour to the action pattern. A list of available behaviour
  will pop up. Every behaviour can be instantiated only once per
  action pattern.
\item Add transition \\
  Click on another action pattern to link the two pattern with a
  transition. A dialog box will appear to enter the transition name.
\item Rename pattern \\
  A dialog box will appear to enter the new name of the action
  pattern.
\item Rename transition \\
  A dialog box will appear with the name and the target of the
  transitions. You can also delete a transition by deleting the name
  or the target of the transition. - Be careful, the new entry
  will only be accepted, if you left the entry field by clicking on
  another entry before hitting the okay button.
\item Delete pattern \\
  Deletes the pattern.
\end{itemize}

Clicking with the right mouse button into the green area of an action
pattern produces a popup menu with the following options.
\begin{itemize}
\item Set arbiter \\
  Select the arbiter for to the action pattern. A list of available arbiters
  will pop up. Every action pattern has to have an arbiter.
\item Delete arbiter \\
  Deletes the arbiter from the action pattern.
\end{itemize}

Clicking with the right mouse button onto a behaviour within an action
pattern produces a popup menu with the following options.
\begin{itemize}
\item Up \\
  Move the behaviour up one position in the list of behaviours. The
  position in the behaviour list represents the priority of the
  behaviour at the arbiter. - The higher, the more important.
\item Down \\
  Move the behaviour down one position in the list of behaviours.
\item Set parameters \\
  The parameters dialog will be displayed. It has three columns. In
  the left one, the parameter names are displayed as label. The middle column
  contains a single line entry field for the parameter. In the right column the
  measure of the parameter is displayed as label, if no measure is available
  the parameters type is displayed. The default value is
  displayed as bubble help associated with each entry field, if
  available. The parameters type is displayed as bubble help
  associated with each measure label. The dialog is strictly
  typed. Therefore the okay button is inactive if any of the entry fields
  are containing invalid input, for instance a value bigger than 180
  for an Angle type.
\item Delete behaviour \\
  Deletes the behaviour from the action pattern.
\end{itemize}

\subsection{Describing the Available Behaviours}
\label{SEC:BehaviourDescription}

To be able to use your own behaviours within the PolicyEditor, you
have to describe their properties (name and parameters) within a
parameters description file. The parameter framework is explained in
more detail in the previous chapter. Anyhow there are some minor
restrictions for behaviour parameter description files, that are
explained below:

$<$!--MiroParametersConfigDocument--$>$

\begin{description}
  \item[config]
  \begin{description}
  \item[config\_group]

    Attributes:
    \begin{itemize}
    \item \textit{name} The name of the config group. Within the
      behaviour parameters framework the valid names are \textbf{behaviour} and
      \textbf{arbiter}. Within each group only config items of the
      specified type are allowed to occur. Other groups are processed
      by the parameter auto code generation tool {\tt MakeParams} but
      ignored by the {\ PolicyEditor}.
    \end{itemize}

    \begin{description}
    \item[config\_item]

      Attributes:
      \begin{itemize}
      \item \textit{name} The name of the behaviour or the arbiter.
      \item \textit{parent} The name of the super class. That has to
        correspond to the corresponding super class of the behaviour.
      \item \textit{namespace}
      \end{itemize}
      \begin{description}
      \item[config\_parameter]

        Attributes:
        \begin{itemize}
        \item \textit{name}
        \item \textit{type} No nested types are allowed!
        \item \textit{default}
        \item \textit{measure}
        \item \textit{inherited}
        \end{itemize}
      \end{description}
    \end{description}
  \end{description}
\end{description}

\subsection{Example Behaviour Description File}

\subsection{Auto-generating Parameter Class Code}
\label{sec:makeParams}

Having written the description of your behaviours parameters within
XML, it is possible to auto-generate the necessary code for your
behaviours associated Parameters class. See the previous chapter for
details.

\section{Policy Controller}
\label{SEC:POLICY_CONTROLLER}

The policy controller is some kind of remote control for the behaviour
engine. It consists of a panel with four buttons and a menu bar. The
buttons allow to start, stop, suspend and resume the behaviour engine.

\subsection{The Menu}

\subsubsection{File}
\begin{itemize}
\item Connect robot \\
  Connect the policy controller with a behaviour
  engine of a robot. A dialog box will be displayed to enter the
  robots name (the naming context, its behaviour engine is registered
  at).
\item Load policy \\
  Load a policy file into the behaviour engine. A dialog box will be
  displayed to enter the file name. The file has to be locally
  accessible by the robots behaviour engine.
\item Send policy \\
  Send a policy to the robot.  The file dialog will be displayed to
  select the policy. The behaviours and arbiters within the policy
  have to be present within the behaviour and arbiter repositories in
  the behaviour engine running on the robot. - Otherwise the parsing
  will fail.
\item Quit \\
  Quit the policy editor.
\end{itemize}

\subsubsection{Edit}
\begin{itemize}
\item Send transition \\
  Send a transition to the behaviour engine. A dialog box will be
  displayed to enter the transition name. This is for debug purposes
  only, as the transition will only succeed if it is registered at the
  currently active action pattern - which in turn can change any time.
  There will be also no feedback whether the transition succeeded or
  not.
\item Send global transition \\
  Send a global transition to the behaviour engine. A dialog box will be
  displayed to enter the action pattern to activate. The action
  pattern has to be present within the current policy, otherwise the
  transition will fail.
\end{itemize}

%%% Local Variables: 
%%% mode: latex
%%% TeX-master: "miro_manual"
%%% End: 

\chapter{Writing a \miro Service}

The programming of a \miro service is divided into two step:
Accessing the low level hardware device (such as the Can bus) and
writing the high level interfaces. Note that many hardware devices
incorporate multiple sensor modalities.


\section{High-Level Server Programming}

\begin{figure}[!ht]
  \ifpdf
  \includegraphics[width=\textwidth]{../fig/client_server.png}
  \else
  \includegraphics[width=\textwidth]{../fig/client_server.eps}
  \fi
  \caption{Client-Server Architecture Overview}
  \label{FIG_CLIENT_SERVER_ARCHITECTURE}
\end{figure}

In order to write a server you have to do the following steps:
\begin{enumerate}
  \item Describe your interface in IDL.
  \item Translate the IDL description to C++ files.
  \item Implement your methods.
\end{enumerate}


\subsection{Copy the \miro server template directory}

- Which files are in there


\subsection{Describing an interface in IDL}

- edit interface description (.idl) \\
- idl/Makefile 

\subsection{Translating the IDL description to C++}

- call TAO\_idl (in {\tt \$(TAO\_ROOT)/TAO\_IDL})

\subsection{Implementing your own methods}

- in ...Impl \\
- don't forget to document !!


\section{Low-Level Server Programming}

Low-Level Servers are Servers that talk to hardware devices. They
usually use a device driver that is compiled into the kernel or loaded 
as a module. They provide handling for special protocols which have to 
be followed and provide the services of the hardware device in an
abstracted way, so that they are useful.

Hardware devices in a UNIX system are represented by file
descriptors. There are two groups, character devices and block devices,
character devices are stream based and provide for sequential reading
and writing, while block based devices provide random access to block
addresses. At the moment all hardware devices used by miro use the
character device abstraction.

A useful \miro low-level server must provide two functionalities:
interaction on the CORBA level with (potential multiple) clients that
use the service concurrently and handling of all communication with
the device, which may be asynchronous to the requests. This requires a
multithreaded design. Fortunately the ACE framework provides design
patterns that fulfill these needs. For a schematic view of a \miro
low-level server see figure
\ref{FIG_CLIENT_SERVER_ARCHITECTURE_LOW_LEVEL}.


\subsection{The Device Framework}

\subsubsection{Connection}

\subsubsection{EventHandler}

\subsubsection{Consumer}

\subsection{The Configuration Framework}

\subsection{Parameter}

\subsection{XML parsing}

\subsection{The Reactor and Events}

The ACE reactor class provides an abstraction of the popular
{\tt select} method in common UNIX programming environments. It is
able to invoke an event handler when data is readable from the file
descriptor\footnote{it can also handle timer events and asynchronous
  writes, see \cite{ACE-Manual} for details}. It provides low
latency and the thread does not consume any CPU power while waiting
for data to come. To keep the responsiveness of this architectural
element high\footnote{which is essential to avoid flooded input
  queues}, only few computation should be done inside of the event
handler, usually this event handler only assembles full packets and
passes these on to a separate task which accounts for higher level
processing. It is important to note that the reactor and the event
handler share a single thread of execution, so don't expect to get
feedback from the reactor within the event handler. 


\subsection{Using Tasks}

Tasks are the second important design pattern provided by ACE that we use.
A task is an assembly of one or more threads, that share a common
message queue, which can be used to pass requests to the threads. We
usually use a task for handling of complete packets assembled by the
event handler, which puts them on the message queue. According to the
type of the packet the tasks takes appropriate action. This usually
consists in retrieving data from the packet and signaling it on a
condition variable. Condition variables are the main mechanism of the
(CORBA) server implementation methods to wait for data requests or
events that are provided by the hardware. The ACE task class provides
an easy scaling for multiple working threads that concurrently work on 
the different requests.

\begin{figure}[!ht]
  \ifpdf
  \includegraphics[width=\textwidth]{../fig/client_server_low_level.png}
  \else
  \includegraphics[width=\textwidth]{../fig/client_server_low_level.eps}
  \fi
  \caption{Client-Server Architecture}
  \label{FIG_CLIENT_SERVER_ARCHITECTURE_LOW_LEVEL}
\end{figure}


\subsection{Thread/Task Synchronization}

\subsection{ACE logging}

\subsection{A Simple Example}

%%% Local Variables: 
%%% mode: latex
%%% TeX-master: "miro_manual.tex"
%%% End: 




\begin{appendix}
\chapter{\miro installation}
\label{sec:miroInstall}

This section is an extended version of the INSTALL file contained
within \miro. If you have problems installing \miro, read on.  Ask
your system administrator if there exists a central installation, or
whether it would be more convenient to prepare a central installation.

\subsection{Requirements}

You need the following packages in order to compile \miro:
\begin{itemize}
\item ACE ($>$= 5.2.2): Please refer to the installation instruction
  in \ref{sec::aceInstall} for building and installing ACE. The source
  tarball and CVS access is available at:\\
  \texttt{http://www.cs.wustl.edu/$\sim$schmidt/ACE.html}.
\item TAO ($>$= 1.2.2): Building and installation is also described in
  section \ref{sec::aceInstall}. TAO downloads are available on:\\
  \texttt{http://www.theaceorb.com/}.
\item Qt ($>$= 2.2.4): Installation should be quite easy, as Qt uses a
  configure script for automatic system checks. Qt is installed on
  most modern Unix/Linux systems anyway, since KDE relies on
  it. Recent versions of Qt can be found at:\\
  \texttt{http://www.trolltech.com/}.
\end{itemize}

If you want to compile the wrapper classes for the speech detection
system, you also need:
\begin{itemize}
\item Speech tools: The Edinburgh Speech Tools Library is a collection
  of C++ class, functions and related programs for manipulating the
  sorts of objects used in speech processing. It is to provide the
  underlying classes in the Festival Speech System. We use version
  1.2.3 that can be retrieved from:\\
  \texttt{http://www.cstr.ed.ac.uk/projects/speech\_tools}
\item Festival: The Festival Speech Synthesis System is a general
  multi-lingual speech synthesis system developed at CSTR. We
  successfully used version 1.4.2. See:\\
  \texttt{http://www.cstr.ed.ac.uk/projects/festival/}
\item Sphinx: The CMU Sphinx Group Open Source Speech Recognition, a
  real-time, large vocabulary, speaker independent speech recognition
  system. We use version Sphinx2-0.4, downloaded from:\\
 \texttt{http://www.speech.cs.cmu.edu/sphinx}
\end{itemize}

Depending on your camera system, you may need libraries for the
IEEE 1394 (aka Firewire) support:
\begin{itemize} 
\item libraw1394: This library provides direct access to the IEEE 1394
  bus through the Linux 1394 subsystem's raw1394 user space
  interface. It can be found one it's homepage:\\
  \texttt{http://www.linux1394.org/}\\
  or on SourceForge:\\
  \texttt{http://sourceforge.net/projects/libraw1394/}
\item libdc1394: It is a library that is intended to provide a high
  level programming interface for application developers who wish to
  control IEEE 1394 based cameras that conform to the 1394-based
  Digital Camera Specification. Available via:\\
  \texttt{http://sourceforge.net/projects/libdc1394/}
\end{itemize}
Both libraries are available as compiled packages for the most modern
Linux systems too.

To build the documentation, your system should provide the following
tools:
\begin{itemize}
\item Doxygen: It is a JavaDoc like documentation system for C++, C,
  Java and IDL. It can generate an on-line HTML documentation
  extracted directly from the sources, which makes it much easier to
  keep the documentation consistent with the source code. You can
  download from the homepage:
  \texttt{http://www.stack.nl/\~dimitri/doxygen/}\\
  or again from SourceForge:\\
  \texttt{http://sourceforge.net/projects/doxygen/}
\item The LaTeX package (either pdflatex or latex and dvips, bibtex
  and makeindex -- available via one of the CTAN server, e.g.\ 
  \texttt{http://dante.ctan.org}) and the convert image convertsion
  utility from the ImageMagick package \\
  (\texttt{http://imagemagick.sourceforge.net/}).
\end{itemize}

\subsection{Download}

As if you read this text, you already have a version of \miro, but if
you want to update this version, or need some additional information,
please have a look at
\texttt{http://smart.informatik.uni-ulm.de/MIRO/index.html}

\subsection{Compilation}

\miro is now shipped with a configure script, so the package can be
installed with the usual steps:

\begin{lstlisting}[frame=tb]{}
./configure
make 
(make install)
\end{lstlisting}

The configure script looks automatically for the necessary and
optional software packages, and will stop or print a warning if some
of them are not found. If you install e.g.\ ACE in an unusual place,
you can pass the configure script some additional options like for
example \texttt{--with-ACE=<ACE-root-dir>}. For Qt, their are more
flexible command-line parameters, because Qt is often distributed over
several directories (like e.g.\ in a standard Debian installation).
Also some Linux distributions ship a multi-threaded version of the Qt
library, called \texttt{qt-mt}.  You can specify this with
\texttt{--with-Qt-lib=qt-mt}. A complete list of command-line options
are accessible via \texttt{./configure --help}.

Another way of telling the configure script where the packages are
installed is the use of environment variables. For the required
packages, these are \texttt{ACE\_ROOT}, \texttt{TAO\_ROOT} and
\texttt{QTDIR} point to the base directory of each package. Because
ACE and TAO need this environment variables anyway, this is the most
comfortable way to point to the directories. For the speech detection
wrappers classes, \texttt{SPEECH\_TOOLS\_ROOT}, \texttt{SPHINX\_ROOT}
and \texttt{FESTIVAL\_ROOT} are possible. It can be necessary, to let
the \texttt{LD\_LIBRARY\_PATH} point to all the libraries anyway.

If the configure script did not find the different packages anyhow, or
if your system require unusual options for compilation or linking that
the configure script does not know about, you can give configure
initial values for variables by setting them in the environment. Using
a Bourne-compatible shell, you can do that on the command line like
this:

\begin{lstlisting}[frame=tb]{}
CFLAGS=-O2 LIBS=-lposix ./configure
\end{lstlisting}

Or on systems that have the `env' program, you can do it like this:

\begin{lstlisting}[frame=tb]{}
env CPPFLAGS=-I/usr/local/include LDFLAGS=-s ./configure
\end{lstlisting}

On runtime, some applications rely on two environment variables,
namely \texttt{MIRO\_ROOT} (the base directory) and \texttt{MIRO\_LOG}
(directory, where logging data is stored), so please set these
variables to the appropriate values too.

Additionally, you can enable or disable several features, like the
support for different robot platforms and for different video devices.

Supported robot platforms:
\begin{itemize}
\item The B21 port is the oldest one and should therefore be the most
  stable, regarding the interfaces. All of the robots hardware
  components are supported by \miro.
\item The Sparrow99 and Sparrow2003 robots are our soccer robots. They
  are our testbed for multi-robot programming. Therefore most new
  technology is first tested on this platform. So even as you do not
  have a Sparrow99 robot (we built them ourselfs), it might be
  interesting to look at the sources for this robot if you are looking
  for group communication technology etc.
\item The Pioneer1 platform is a complete port for the ActiveMedia
  robot series, based on the PSOS, P2OS and AROS protocols. But we
  could only test and extend it to the platforms available to us.
  \begin{itemize}
  \item The old Pioneer1 is supported, but lacks maintenance, as our
    model is not used at the lab anymore.
  \item The Performance PeopleBot is in active development and already
    mostly complete. Motors, Sonars, Bumpers and Video are working.
    PanTilt and Gripper are experimental, Zoom is to come next.
  \end{itemize}
\end{itemize}

Supported video devices:
\begin{itemize}
\item Bttv Frame Grabbers: These frame grabber cards are supported via
  video for Linux.  This is the standard way of connecting standard
  analog video cameras to the computer.
\item Firewire Digital Cameras: \miro supports the fire wire digital
  camera protocol using libraw1394 and libdc1394.
\item Matrox Meteor Frame Grabbers: These rather old frame grabber
  cards are also supported by \miro. This device however, is mostly
  unmaintained.
\end{itemize}

Finally, you can choose if the documentation should be build or not.
Even if you decide not to build it, but the configure script found all
the necessary tools, the Makefiles are prepared. So you can go to the
\texttt{doc/tex} and \texttt{doc/html} directory later on and build
the documentation there with a simple make.

After running the configure script, \miro show up a summary of what
will be compiled and which features will not. If this is not what you
desired, please check the messages coming up during the configure run
for packages, \miro did not find.

\subsection{Installation}

\miro can be installed with the a simple:

\begin{lstlisting}[frame=tb]{}
make install
\end{lstlisting}

By default, the package's files will be installed in \texttt{/usr/local/bin},
\texttt{/usr/local/lib} etc.  You can specify an installation prefix other
than \texttt{/usr/local} by giving configure the option
\texttt{--prefix=PATH} (note: if you use the prefix-option, you should
install \miro actually, otherwise libtool may look for the library in
the installation directory and therefore do not find it).

Beside that, \miro can be used already without installation. Therefore,
a make run install the libraries and the binaries during the
compilation into the \texttt{lib/} and \texttt{bin/} directory. This
enlarges the \miro directory, but can be quite useful, if you work on
\miro itself and a derived application at the same time.

Moreover with this trick it is possible, to use different \miro
versions in parallel at the same time. Therefore the build process
allow to compile the package in a different directory. Assuming you
have extracted \miro in a directory called \texttt{Miro}, you can
build it in two different directories with miscellaneous options:

\begin{lstlisting}[frame=tb]{}
mkdir Miro.B21
cd Miro.B21
../Miro/configure --enable-B21
make
cd ..
mkdir Miro.Sparrow99
cd Miro.Sparrow99
../Miro/configure --enable-Sparrow99
make
\end{lstlisting}

\subsection{Additional make targets}

Beside the already described \texttt{install} target and an solely
\texttt{make} (which in fact means \texttt{make all}), there are a
couple of other targets, that may be of interest:

\begin{description}
\item[clean:] Removes object files, libraries, binaries and
  automatically generated code, like the parameter classes generated
  from the xml files and the client and server stubs generated from
  the idl files.
\item[distclean:] Removes all of the above files and the automatically
  generated \texttt{Makfiles} generated from the autotools. So you
  have to reconfigure the packet after this.
\item[dist:] Generates a gzipped tar-file containing a snapshot of the
  actual \miro directory, including all the necessary source files to
  configure and build \miro.
\item[distcheck:] Some as above plus a complete make of the packed
  tar.gz file and finally make another tarfile to ensure the
  distribution is self-contained. You can change the configure options
  that are used for this make run within the toplevel
  \texttt{Makefile.am} by changing the
  \texttt{DISTCHECK\_CONFIGURE\_FLAGS} variable.
\end{description}

\subsection{Developer information}

The build-process within \miro is done using automake, autoconf and
libtool. Normally, there is no need to use these programs, as of
running the configure script is the only thing a user have to do.

The configure shell script attempts to guess correct values for
various system-dependent variables used during compilation. It uses
those values to create a Makefile in each directory of the package.
It also create one config.h file containing system-dependent
definitions that can be used in the source file to allow appropriate
compilation:

\begin{lstlisting}[frame=tb]{}
#ifdef HAVE_CONFIG_H
#include <config.h>
#endif
...
#ifdef MIRO_HAS_1394
...
#endif
\end{lstlisting}

If you want to change something within \miro like for example adding
additional files or directories, you have to change the proper
\texttt{Makefile.am}. Do \textbf{not} change the \texttt{Makefile.in}
or even the final \texttt{Makefile}, because changes there are getting
lost at the next automake run. The same applies, if you want to
change the behaviour of the configure script itself, for example to
detect new software packages or special version information, you have
to change the \texttt{configure.ac} and \textbf{not} the
\texttt{configure} file itself.

Assuming that you have written appropriate \texttt{Makefile.am} and
\texttt{configure.ac} files, you should be able to build your project
by entering the following commands:

\begin{lstlisting}[frame=tb]{}
aclocal -I config
autoheader
automake
autoconf
./configure
\end{lstlisting}

For every consecutive step, you can simply call the resulting
\texttt{config.status} script, that rerun the configure procedure with
the used configuration values.

Depending on your versions of automake, autoconf and libtool you may
get some warning messages about improper usage of \texttt{LDFLAGS}.
Ignore this, since it is without any influence on the final result. We
tested the whole procedure successfully with at least the following
versions:
\begin{itemize}
\item automake 1.5
\item autoconf 2.52
\item libtool 1.4.2
\end{itemize}

Normally, it should be enough, to use the available examples within
\miro (see section \ref{sec:makefile}) to add some small additions to
the build process, but sometimes a deeper understanding of the used
tools might be necessary. Please have a look at the manuals, available
e.g.\ directly from the GNU homepage:
\texttt{http://www.gnu.org/manual/} or read the freely available book
on \texttt{http://sources.redhat.com/autobook/}.

\chapter{ACE and TAO Installation}
\label{sec::aceInstall}

ACE/TAO is a large software package with many configuration
options affecting the build process as well as its runtime
features. Therefore we summarize here our experience with the
ACE/TAO installation.

At the time of this writing we suggest to use the latest BFO (bug fix
only) version ACE/TAO --- currently this is 5.3.2/1.3.2. There are
knon problems with the notification service in TAO.

The options that should be set additionally in the {\tt
  platform\_macros.GNU} file (to be found under
$<$ACE-directory$>$/include/makeinclude/) are the following:

\begin{lstlisting}[frame=tb]{}
debug=0
\end{lstlisting}

Debug information has a huge impact on the footprint of the libraries,
so as long as you do not really need it, disable it.

Additionally, for further footprint reduction you
might also add the following:

\begin{lstlisting}[frame=tb]{}
qt_reactor=0
xt_reactor=0
TAO_ORBSVCS= Naming Notify
DEFFLAGS= -DACE_USE_RCSID=0
ACE_COMPONENTS= FOR_TAO
\end{lstlisting}

The two reactors are not needed by ACE/TAO/Miro. So as long as you do
not use them yourself, they might just be left out. Also Miro only
uses the CORBA Naming Service and the CORBA Notification Service.
Therefore the other services of TAO need not to be built.

We also recommend to read the installation instructions provided
by ACE carefully. We admit, the installation isn't entirely trivial.

To compile and run programs that use the ACE/TAO toolkit, also the
environment variables {\tt ACE\_ROOT} and {\tt TAO\_ROOT} need
to be set to the appropriate root directories of your ACE/TAO
installation. For the bash shell the following lines in your local
{\tt .bashrc} file should do the job:

\begin{lstlisting}[frame=tb]{}
export ACE_ROOT=<path to the ACE directory>
export TAO_ROOT=$ACE_ROOT/TAO
export LD_LIBRARY_PATH=$ACE_ROOT/ace/:$LD_LIBRARY_PATH
\end{lstlisting}
%%$

%%% Local Variables: 
%%% mode: latex
%%% TeX-master: "miro_manual"
%%% End: 

\chapter{Project setup with Automake and Autoconf}
\label{sec:new_project}

This chapter tries to give a short introduction on how to set up your
own project that uses \miro by the help of automake and autoconf.
This text was created using autoconf 2.52 and automake 1.5, so your
local installation may differ slithly.

First describe the the mode of operation of automake and autoconf is
explained by a real-world example from the \miro source
base. Afterwards the creation of the root configuration file of the
autoconf tool, {\tt configure.ac} is discussed. Afterwards the
different Makefile.ac templates, shipped with \miro are
introduced. After a short survey on the actual buid directives, the
chapter concludes with links to internet resources on the auto-tools.

\section{Introductory Example}

First a short introduction into the whole structure of automake and
autoconf, that uses a two-stage process.

\begin{itemize}
\item With automake you can (normaly) easily describe, how the files
within a folder should be processed. The commands for this are written
in a file called Makefile.am. As a minimal example look for example in
{\tt \$MIRO\_ROOT/\-src/\-can/\-Makefile.am}:

\begin{verbatim}
  lib_LIBRARIES = libcan.a
  libcan_a_SOURCES = Parameters.xml CanConnection.cpp CanMessage.cpp ...
\end{verbatim}

This generates a (static) library named {\tt libcan.a}, which is build
from the assigned sources {\tt Parameters.xml, CanConnection}...  All
other stuff within the Makefile.am is e.g.\ which files are installed
into which directory, if the user applies a {\em "`make install"'\/}
({\tt can\_include\_HEADERS}, note the prefix of the macro), or which
files are generated automatically during the build-process and can
therefor removed if a {\em "`make clean"'\/} is desired ({\tt
BUILT\_SOURCES} and {\tt CLEANFILES}). These Makefile.am are converted
into {\tt Makefile.ac} using the {\em automake\/} program.

\item {\em autoconf\/} generates the real {\tt Makefile} for the {\tt
Makefile.ac} files. Therefor it has to know the paths and options
necessary for the build process and which files, libraries and
programs it has to compile at all. Exactly that is, what the
configure-script tries to detect and guess automatically. Of course no
one wants to write a configure-script directly (take a look at it, it
looks really, really ugly. Moreover a program should compile on a
couple of different systems, each having their own quirks and
peculiarities). Therefor one writes a {\tt configure.ac} files, that
is finally converted into the configure-script. The whole autoconf
stuff is based on so called macros, from which are a lot of them
allready included:

\begin{verbatim}
  AC_CHECK_LIB(raw1394, raw1394_get_libversion, ac_have_libraw1394=yes, [ac_have_libraw1394=no])
  AC_CHECK_HEADERS(errno.h fcntl.h stdlib.h)
  AC_CHECK_PROG(ac_has_dvips, dvips, yes, no)
\end{verbatim}

There are e.g.\ macros to test the existence of libraries, header files
or executable programs. If you need a more complicated test, you can
write them with a language called M4 (M4's syntax is quite close to a
normal shell syntax). The results of all tests are write in a file
called config.h additionally. This file can be included into the
C-files and are treated as normal c-defines (e.g. {\tt \#define
MIRO\_HAS\_METEOR 1}).
\end{itemize}


\section{Create {\tt configure.ac}}

The easiest way to start with is using the autoscan program (from the
autoconf-package). This will create a file called configure.scan that
you should rename to configure.ac. ({\em autoscan\/} is also a good
tool to check your existing configure.ac for missing tests). Now look
into your new configure.ac file. The first line ({\tt AC\_INIT})
describe your project, so fill it with useful values, a name for your
package, your version and an email address for bug-reports. The next
two lines ({\tt AC\_CONFIG\_SCRDIR}, {\tt AC\_CONFIG\_HEADER}) should be
fine.

The remainder of the file is group into different parts. First part is
the check for programs. {\em autoconf\/} offers several ready to use checks
for this, like e.g. {\tt AC\_PROG\_CC} that checks several behaviours of
your installed C-compiler or {\tt AC\_PROG\_LN\_S }that looks how to
create a symlink. The automatically included macros should be fine.

The next paragraph checks for libraries. Note, that the philosophy of
{\em autoconf\/} is to check for the possibility of creating and
compiling a program with a certain function using this library, not
the pure existence of this library. So look for a function that should
be within the different libraries and insert them in the second field
of the {\tt AC\_CHECK\_LIB} macros (the first name is the library name
itself). Possibly you have to remove lines for libraries, that are in
fact builded by your project itself (Note, that it may be difficult to
use C++ specific stuff here).

Third, header files are checked. Beside some prefabed macros for
standard headers (e.g. {\tt AC\_HEADER\_STC}) you can insert additional
files into the {\tt AC\_CHECK\_HEADERS} macros. Insert them as a space
separated list.

Next two paragraph contains checks for typedef and structures and
library functions. The automatically included macros should be fine,
otherwise include additional functions into the {\tt AC\_CHECK\_FUNCS}
macro. The functions here refer to standart C functionality, not to
mix up with special functions provided by additional libraries as
described above.

The last, long macro {\tt AC\_CONFIG\_FILES} contains a space separated
list of all the Makefiles the configure script has to produce. This
means, if you add a directory into your project later on, you have to
insert the path here.



\subsection{Tests for more complex packages (and facilities)}

Most checks for the functionallity of your system should be covered
with the above, autogenerated checks. But to build a project that
relies on other, more complex packages (in our case \miro, ACE, TAO,
Qt, etc), it is necessary to add additionally tests to the configure
script.

Additionally checks are written in a language called M4. It is up to
you, where to save these files, but I prefer to add an extra directory
named config, so the following explanation will assume this.

First you have to check for the existence of e.g. ACE, TAO and
Qt. Simply copy the necessary files {\tt search\_ace.m4},{\tt
search\_tao.m4} and so on from {\tt \$MIRO\_ROOT/\-config/} to your local
config directory and add {\tt AC\_SEARCH\_ACE}, {\tt AC\_SEARCH\_TAO} and
so on (the name of the function macros) to the {\tt configure.ac}
script (for example directly above the {\tt AC\_CONFIG\_FILES} macro.

To be sure to add all the necessary libraries and options that are
used to compile ACE/\-TAO, additional test are prepared within the file
{\tt feature\_ace.m4} (again copy this file into your local folder and
add the macro name into {\tt configure.ac}).

To check for Miro and parts of the functionality, we prepared two
files {\tt \$MIRO\_ROOT/\-templates/\-search\_miro.m4} and {\tt
\$MIRO\_ROOT/\-templates/\-features\_miro.m4}. Again use them the usual way,
e.g. insert {\tt AC\_FEATURES\_BTTV\_MIRO} to be sure that Miro is
compiled with support for {\em BTTV\/} frame grabber cards.

If you need to write your own tests and checks, please feel free to
use and modify these files. Additional macros that can be used out of
the box or as examples for own tests, look at:

\begin{verbatim}
http://wwww.gnu.org/software/ac-archive/
http://ac-archive.sourceforge.net/
\end{verbatim}

\section{Create {\tt Makefile.am} for directories, libraries and executables}

The next step is to prepare a Makefile.am within each directory that
contain files that need to be compiled. The easisest way is again to
use the examples within {\tt \$MIRO\_ROOT/\-templates/}. For example to
create a Makefile.am that should call other Makefiles in other
subdirs, take {\tt \$MIRO\_ROOT/\-templates/\-Makefile.am.dir}, rename it to
{\tt Makefile.am} and add the names of the subdirectories. There are also
examples for sources that should be compiled into a static library
(Makefile.am.lib.a) into a shared library (Makefile.am.lib.so) or
different executables (Makefile.am.bin). All examples contain some
comments to guide you, what to fill in.

Remeber to be sure, that each directory with a {\tt Makefile.am} is
listed in the {\tt AC\_CONFIG\_FILES} macro in the {\tt configure.ac}
script.


\subsection{Makefile.am.dir}
This Makefile starts subbuilds in the specified directories. It is an
example for conditional compilation (depending on the
configuration). If you don't have conditional branches, you can left
the second variable empty or reduce the Makefile.am to a single line,
like e.g. \texttt{SUBDIRS = firstDir secondDir}.

\subsection{Makefile.am.bin}
This Makefile builds a single or several binaries. The name of the
binaries to be built are specified in the variable {\tt bin\_PROGRAMS}.
There have to be an extra line for each binary to be build starting
with the name of the binary itself and an additional
\texttt{\_SOURCES}-suffix. The sources are specified with their suffix
(.cpp). Please also list all the header files here, otherwise they
don't get distributed by an call to \texttt{make dist}. If the
programs need to be linked against several libraries, there is again
one line per binary containing the name of the binary plus an
additional \texttt{\_LDADD}-suffix.

\subsection{Makefile.am.qt}
This is a special Makefile.am template for Qt programs. For
simplicity, the Makefile is designed for a single binary per
directory, but beside that, most of the above section also applies
here. The only exception is, that we have to distinguish between
normal source files (variable \texttt{sources}) and files (variable
\texttt{tomocsources}) that need to be compiled with Qt's Meta Object
Compile (\texttt{moc}). Be sure to \textbf{not} include the header
files into the \texttt{tomocsources} variable, but to the
\texttt{sources} variable, otherwise they get "cleaned".

\subsection{Makefile.am.lib.a}
This Makefile builds a static library. The name of the library can be
specified in the variable \texttt{lib\_LIBRARIES}, the files building
the library in the variable \texttt{sources} (with the .cpp suffix; or
.xml in case of an xml-file for \miro's parameter framework). Headers
that are not automatically considered because they don't have a
corresponding cpp-file should be specified in the variable
\texttt{extraheader}. Don't forget to name the prefix for the
\texttt{\_SOURCES}-macro with the name of the resulting library.
Therefore please substitued all occuring dots with underscores (e.g.\ 
the usual \texttt{.a} a the end of an static library).  Additionally,
replace the prefix for the macro \texttt{\_include\_HEADERS} with the
name of the directory, to which the header files should be installed
(see also the description of \texttt{Make-rules} below).

\subsection{Makefile.am.lib.so}
Makefiles for shared libraries look quite the same as for static
libraries, to most of the above applies here too. The only difference
from a users perspective is the name of the variable
\texttt{lib\_LTLIBRARIES} where you name your library (note the
\texttt{LT} within the variable (and the targets at the end of the
\texttt{Makefile.am}) which indicates that the library is build using
the libtool package).

\subsection{Conditionally Compiled Sub-projects}

If you programs or libraries need to be compiled differently depending
on configuration options, please have a look at an existing
\texttt{Makefile.am}, e.g.\ in
\texttt{\$MIRO\_ROOT/\-examples/\-behaviours/\-engine/} or
\texttt{\$MIRO\_ROOT/\-src/\-miro/}.

\subsection{Create other {\tt Makefile.am}}

Because auto-conf/\-make themself are quite unflexible if you intend to
do something, that the autotools are not designed to, you can also
insert an "`all-local"' target into the {\tt Makefile.am} that is
treated as a normal Makefile part. Although this has to be necessarry
sometimes, be careful with this, auto-conf/make may fail to calculate
the correct dependencies which in turn may lead to incomplete and
corrupt builds. Save places to write your own stuff are for example in
the documentation directory. See {\tt \$MIRO\_ROOT/\-doc/\-html} and {\tt
\$MIRO\_ROOT/\-doc/\-tex} for examples on building a doxygen or latex
documentation.

\subsection{Make-rules}
Last you may have a look at the \texttt{Make-rules} file within the
\texttt{\$MIRO\_ROOT/\-config/} directory. This file is included into each
\texttt{Makefile.am}, so rules or variables that are needed globally,
can be inserted here. If you want to define additional installation
subdirectories for the header files, you can see examples here.

\section{Build the beast}

To finally build the new project, you have to call the different
programs, which in turn convert the prepared files into a real,
make-based project.

Assuming that you have written appropriate {\tt Makefile.am} and {\tt
configure.ac} files, you should be able to build your project by
entering the following commands:

\begin{description}
\item[autoheader] Creates \texttt{config.h.in}.
\item[aclocal -I config] Adds \texttt{aclocal.m4} to directory.
  Defines some m4 macros used by the auto tools and the self-written
  tests from the config-directory.

\item[autoconf] Creates \texttt{configure} from \texttt{configure.ac}
\item[automake] Creates \texttt{Makefile.ac} from \texttt{Makefile.am} 
\item[./\-configure] Creates \texttt{Makefile} from \texttt{Makefile.in}
\item[make]
\end{description}

Just repeat the last 5 steps to completely rebuild the project. Most
projects have an \texttt{autogen.sh} script that runs everything up to
the configure step.


\section{Additional sources of help}

In case of problems, first have a look at the offical manuals of both,
automake and autoconf. You can find them on the GNU webpage:

\begin{verbatim}
http://www.gnu.org/software/automake/manual/automake.html
http://www.gnu.org/software/autoconf/manual/autoconf-2.57/autoconf.html
\end{verbatim}

There is also a book available under the terms of the Open Publication
License:

\begin{verbatim}
http://sources.redhat.com/autobook/download.html
\end{verbatim}

%%% Local Variables: 
%%% mode: latex
%%% TeX-master: "miro_manual"
%%% End: 

\end{appendix}
%------------------------------------------------------------------------------


\bibliography{miro_manual}
\printindex

\end{document}
